\documentclass[stage3a]{tnreport}
% \documentclass[confidential]{tnreport} % À activer si ton mémoire chez Wavestone est confidentiel

% ==============================================================================
% METADONNÉES DU MÉMOIRE (IAE & TELECOM Nancy)
% ==============================================================================
\def\reportTitle{L'intégration de l'IA générative dans le conseil} 
\def\reportLongTitle{L'intégration des outils d'intelligence artificielle générative dans le conseil : comment l'automatisation des tâches d'analyse redéfinit-elle la légitimité et le développement des compétences des consultants ?} 

\def\reportAuthor{Arthur SAUNIER}
\def\reportAuthorEmail{\email{arthur.saunier@telecomnancy.eu}} % Ton mail étudiant

% Configuration de l'adresse (Optionnel, à remplir ou laisser vide)
\def\reportAuthorAddress{} 
\def\reportAuthorCity{} 
\def\reportAuthorPhone{} 

% Encadrement
\def\reportSupervisor{M. Riccardo BERGAMASCHI} % Tuteur professionnel Wavestone
\def\reportAcademicSupervisor{M. Ferri BRIQUET} % Tuteur académique IAE Nancy

% Informations Entreprise
\def\reportCompany{Wavestone Luxembourg} 
\def\reportCompanyAddress{}  
\def\reportCompanyCity{Luxembourg} 
\def\reportCompanyPhone{} 
\def\reportCompanyLogoPath{figures/company-logo} % Ton logo Wavestone dans le dossier figures

\def\place{Nancy} % Ville pour l'engagement anti-plagiat
\def\date{\today} 

% Ligne de signatures réglementaire de fin de document
\addSignaturesLine{{
  {figures/school-logo}{Arthur SAUNIER}{Élève-Ingénieur \& Étudiant MAE},
  {figures/company-logo}{Riccardo BERGAMASCHI}{Tuteur Professionnel}
}}
\addSignaturesLine{{
  {figures/school-logo}{Ferri BRIQUET}{Tuteur Pédagogique IAE}
}}

% ==============================================================================
% DÉBUT DU DOCUMENT
% ==============================================================================
\begin{document}

\maketitle
\pagenumbering{roman}

% Engagement anti-plagiat (Remplacer avec ton numéro étudiant et ta signature){figures/ta_signature}
\insertAntiPlagiarismAgreement{Saunier, Arthur}{eouorjier}{0.2}
\cleardoublepage{}

%\makesecondtitle{}

% ==============================================================================
% REMERCIEMENTS
% ==============================================================================
\section*{Remerciements}
\addcontentsline{toc}{chapter}{Remerciements}

Je tiens à remercier tout particulièrement... % À compléter plus tard

\cleardoublepage{}

% ==============================================================================
% TABLE DES MATIÈRES (Sommaire synthétique obligatoire pour l'IAE)
% ==============================================================================
\renewcommand{\baselinestretch}{0.5}\normalsize
\tableofcontents
\renewcommand{\baselinestretch}{1.0}\normalsize
\cleardoublepage{}

\pagenumbering{arabic}
\setcounter{page}{1}

% ==============================================================================
% INTRODUCTION GENERALE (Validée ensemble)
% ==============================================================================
\chapter{Introduction Générale}

\section{Contexte sectoriel : Le conseil à l'ère de l'intelligence artificielle}
Le secteur du conseil en stratégie et en management traverse une période de mutation technologique sans précédent. Historiquement définies comme des \textit{Knowledge-Intensive Firms} (KIFs), les entreprises de conseil tirent leur valeur ajoutée de la capacité de leurs collaborateurs à collecter, analyser et synthétiser des informations complexes pour éclairer la prise de décision des dirigeants. 

L'émergence et la démocratisation fulgurante des outils d'intelligence artificielle générative depuis la fin de l'année 2022 sont venues bousculer ce modèle économique et opérationnel. Des tâches autrefois dévolues aux consultants juniors, telles que la revue documentaire, la synthèse de rapports sectoriels ou la première formulation de recommandations, peuvent désormais être automatisées ou fortement accélérées par des modèles de langage de grande taille (LLM). 

Cette transformation technologique offre des opportunités majeures en termes d'efficience opérationnelle et de gain de productivité pour les cabinets. Toutefois, elle soulève des questionnements managériaux profonds. D'une part, l'automatisation des tâches d'analyse de bas niveau risque d'altérer le processus traditionnel d'apprentissage et de développement des compétences des collaborateurs en début de carrière. D'autre part, la récurrence des biais algorithmiques et des hallucinations informationnelles pose un risque de réputation et de responsabilité juridique majeur pour des structures dont la "confiance" constitue le principal actif immatériel.

\section{Présentation du cadre de l'étude : Wavestone Luxembourg}
C'est au cœur de cette dynamique que s'inscrit ce travail de recherche, mené dans le cadre d'un stage de fin d'études au sein du cabinet Wavestone Luxembourg. Acteur majeur du conseil en transformation digitale et en cybersécurité, Wavestone constitue un terrain d'observation privilégié pour analyser la dualité de l'intégration de l'IA. En tant que consultant au sein de la practice Cybersécurité, confronté quotidiennement à l'intensité des missions de gestion de crise et d'évaluation de la maturité des organisations, l'opportunité nous est donnée d'observer directement la manière dont les équipes s'approprient — ou rejettent — ces nouveaux outils transformationnels.

\section{Énoncé de la problématique}
Face à ce constat, l'objectif de ce mémoire professionnel n'est pas d'analyser le processus technique de développement des outils d'IA, mais bien d'en étudier les impacts organisationnels et managériaux[cite: 2]. Il s'agit de comprendre comment l'organisation navigue entre l'impératif stratégique d'innovation et la préservation de son capital humain et de sa fiabilité. 

Dès lors, la réflexion sera guidée par la problématique suivante : \textbf{L'intégration des outils d'intelligence artificielle générative dans le conseil : comment l'automatisation des tâches d'analyse redéfinit-elle la légitimité et le développement des compétences des consultants ?}

\cleardoublepage{}

% ==============================================================================
% PARTIE 1 : CADRE THÉORIQUE
% ==============================================================================
\part{Cadre Théorique -- L’Expertise face à l’Automatisation}

\chapter{Les Sociétés de Conseil (KIFs) et le modèle de l'expert}
\section{La structure des organisations intensives en connaissances}

Le concept de firme intensive en connaissances (\textit{Knowledge-Intensive Firms} ou KIFs) est essentiel pour appréhender la dynamique organisationnelle d'un cabinet de conseil tel que Wavestone. Contrairement aux entreprises industrielles classiques où le capital est principalement matériel (machines, usines), les sociétés de conseil reposent sur un capital presque exclusivement immatériel : le savoir, l'expertise et la capacité d'analyse de leurs collaborateurs.

Selon la typologie d'Henry Mintzberg (1982), le cabinet de conseil s'apparente à une « bureaucratie professionnelle ». Dans ce modèle organisationnel, le centre opérationnel est composé de professionnels hautement qualifiés qui disposent d'une large autonomie dans leur travail. La coordination ne se fait pas par une supervision hiérarchique stricte ou par une standardisation des processus de production, mais par la standardisation des qualifications. En d'autres termes, l'organisation fait confiance à l'expertise accumulée par le consultant pour répondre de manière sur-mesure aux problématiques complexes de ses clients.

Toutefois, cette autonomie et cette légitimité ne sont pas innées ; elles se construisent et se maintiennent à travers une structure sociale et des mécanismes de réputation rigoureux, où la confiance du client constitue le pilier central de la relation d'affaires.

\subsection{Le processus d'apprentissage et de construction de l'expertise chez le junior}

Au sein de ces structures, l'intégration et la formation des consultants juniors répondent à une logique de compagnonnage industriel et d'apprentissage par la pratique (\textit{learning by doing}). Le consultant débutant est traditionnellement affecté à des tâches dites « de base » mais fondamentales : 
\begin{itemize}
    \item La recherche documentaire et la veille réglementaire.
    \item La collecte et le traitement de données brutes sur le terrain.
    \item La première synthèse de rapports ou la rédaction de livrables intermédiaires.
\end{itemize}

Ces activités, bien que parfois perçues comme chronophages ou à faible valeur ajoutée immédiate, remplissent une fonction pédagogique critique. Elles permettent au junior de s'imprégner de la culture technique du domaine (comme les frameworks de gestion de crise cyber ou les normes de l'ENISA), de développer son esprit critique et de légitimer progressivement sa posture d'expert. C'est en se confrontant à la complexité de la donnée brute que le cerveau humain apprend à structurer un raisonnement managérial.

L'introduction d'outils capables d'automatiser instantanément ces phases de recherche et de synthèse vient donc directement percuter ce modèle traditionnel de transmission du savoir. Si le consultant junior délègue la phase d'assimilation et de traitement critique de l'information à une intelligence artificielle, la question centrale devient celle de sa capacité future à développer une expertise réelle, autonome et légitime. % Idéalement, crée un sous-dossier "chapters"

\chapter{L'IA Générative comme technologie de rupture organisationnelle}
\section{La firme intensive en connaissances et son modèle d'apprentissage}

Le chapitre précédent a décrit une situation de travail~; il reste à la doter des
instruments qui permettent de l'analyser. Trois corps de littérature sont mobilisés
successivement. Le premier caractérise ce qu'est structurellement un cabinet de conseil et
comment s'y forme un professionnel. Le deuxième explique sur quoi repose l'autorité de ce
professionnel devant son client. Le troisième examine ce que l'introduction d'une
technologie générative déplace dans cet édifice, en s'appuyant sur les travaux empiriques
les plus proches du terrain étudié. L'ordre n'est pas indifférent~: la question de
l'intelligence artificielle n'est traitée qu'en dernier, une fois établi ce qu'elle vient
percuter.

\subsection{Ce que produit une organisation dont le capital est le savoir}

Le concept de firme intensive en connaissances (\textit{knowledge-intensive firm}) désigne
les organisations dont l'activité repose sur un travail intellectuel complexe conduit par
une main-d'{\oe}uvre très qualifiée, et dont le capital réside dans l'expertise de ses
membres bien plus que dans ses actifs matériels \cite{starbuck1992}. Contrairement à
l'entreprise industrielle, où le capital est constitué de machines et d'installations, le
cabinet de conseil ne possède que ce que ses collaborateurs savent faire. Ce capital sort
de l'immeuble chaque soir, ce qui suffit à expliquer l'attention portée par ces firmes à la
rétention, à la formation et à la codification interne des savoirs.

La taxonomie des \textit{professional service firms} proposée par von Nordenflycht
\cite{nordenflycht2010} précise cette caractérisation par trois traits distinctifs~:
l'intensité de savoir, la faible intensité capitalistique et le degré de
professionnalisation de la main-d'{\oe}uvre. Le conseil en management occupe dans cette
taxonomie une position particulière. Il partage avec la médecine ou le droit l'intensité de
savoir et la faible intensité capitalistique, mais il ne dispose ni d'un ordre
professionnel, ni d'un diplôme d'exercice obligatoire, ni d'un code déontologique
opposable. Cette position intermédiaire a une conséquence directe sur le raisonnement
conduit dans ce mémoire~: faute de barrière institutionnelle à l'entrée, la légitimité du
consultant se rejoue à chaque mission plutôt qu'elle ne se déduit d'un titre. Elle autorise
par ailleurs, avec prudence, à transposer au conseil des résultats obtenus sur d'autres
professions intellectuelles, à condition de tenir compte de cet écart de
professionnalisation.

Alvesson \cite{alvesson2004} spécifie l'analyse au cas des cabinets et introduit la notion
décisive pour la suite~: l'ambiguïté de l'\textit{output}. Le livrable d'une mission de
conseil n'est pas mesurable comme l'est une pièce usinée. Sa qualité ne se constate pas de
façon indiscutable, ni au moment de la remise, ni souvent après. Le client ne peut
généralement pas vérifier par lui-même la justesse d'un diagnostic~; s'il le pouvait, il
n'aurait pas acheté la prestation. Il en résulte que le jugement porté sur le travail de
connaissance emprunte largement à des signaux indirects~: la réputation du cabinet, la
qualité formelle du document, l'assurance de celui qui le présente. Alvesson
\cite{alvesson2001} pousse ce constat plus loin en soutenant que le travail de connaissance
se juge sur l'image et l'identité autant que sur le résultat. Cette proposition, qui paraît
sévère, est en réalité l'articulation centrale du présent chapitre~: si la valeur perçue
tient pour partie à des signaux, alors toute technologie qui modifie la production de ces
signaux — la vitesse de rédaction, la densité apparente d'un document, la couverture
apparente d'un sujet — touche au c{\oe}ur de la prestation et non à sa périphérie.

\subsection{La bureaucratie professionnelle et la standardisation des qualifications}

La typologie des configurations organisationnelles proposée par Mintzberg
\cite{mintzberg1982} fournit la grille structurelle qui manque à ce qui précède. Le cabinet
de conseil s'apparente à une \textit{bureaucratie professionnelle}. Dans cette
configuration, le centre opérationnel est constitué de professionnels hautement qualifiés
qui disposent d'une large autonomie dans la conduite de leur travail. La technostructure y
est faible et la ligne hiérarchique mince, parce que l'essentiel de la coordination ne
passe ni par la supervision directe ni par la standardisation des procédés.
% Pagination à compléter par Arthur sur l'exemplaire consulté :
% \cite[p.~309]{mintzberg1982} pour le chapitre consacré à la bureaucratie professionnelle.

Le mécanisme de coordination propre à cette configuration est la \textit{standardisation des
qualifications}. L'organisation ne prescrit pas comment le travail doit être exécuté~; elle
s'assure en amont que celui qui l'exécute a été formé à le faire correctement. Le contrôle
est déplacé de l'acte vers la personne, et de l'exécution vers la formation. C'est là que le
modèle du conseil se distingue de la bureaucratie professionnelle canonique. Dans la
médecine ou l'enseignement, la standardisation des qualifications est assurée par des
institutions extérieures~: universités, ordres, examens d'État. Dans le conseil, ainsi que
le laissait prévoir la taxonomie de von Nordenflycht \cite{nordenflycht2010}, elle est
largement produite \textit{par le cabinet lui-même}. Le grade, la revue par le manager,
l'exposition graduée à la relation client et le passage par des missions de complexité
croissante remplacent le diplôme d'exercice.

Cette observation a une portée analytique forte. Elle signifie que le dispositif de
formation interne d'un cabinet n'est pas une fonction support parmi d'autres~: il est le
mécanisme de coordination principal de l'organisation. Une perturbation du parcours
d'apprentissage des juniors n'est donc pas un problème de gestion des ressources humaines
isolé, mais une atteinte à ce qui tient la structure ensemble. La question directrice
relative à l'impact de l'automatisation sur la montée en compétences des juniors devient,
sous cette grille, une question de théorie des organisations autant que de GRH.

\subsection{La conversion des connaissances et l'apprentissage par la pratique}

Reste à préciser \textit{comment} cette qualification se forme. La distinction établie par
Polanyi \cite{polanyi1966} entre savoir tacite et savoir explicite en fournit le point de
départ. Le savoir explicite est formulable, transmissible par le langage et stockable dans
un document~: une méthodologie, une norme, un modèle de livrable. Le savoir tacite est
incorporé dans la pratique et résiste à la mise en mots. Il regroupe le jugement, le sens
de la pertinence, la capacité à sentir qu'un chiffre est douteux ou qu'une recommandation
ne passera pas devant un comité. C'est précisément cette part du savoir professionnel qui
distingue un consultant expérimenté d'un débutant disposant des mêmes documents.

Nonaka \cite{nonaka1994}, puis Nonaka et Takeuchi \cite{nonaka1997}, opérationnalisent
cette distinction en un modèle de conversion des connaissances articulant quatre modes~:
socialisation (du tacite vers le tacite), externalisation (du tacite vers l'explicite),
combinaison (de l'explicite vers l'explicite) et internalisation (de l'explicite vers le
tacite). L'intérêt de ce modèle pour le présent travail tient à ce qu'il permet de nommer
précisément ce qui est en jeu. Deux modes seulement sont directement concernés par
l'automatisation de la production documentaire.

La \textit{socialisation} d'abord. Elle se joue dans le partage d'expérience directe~: le
junior qui observe un senior conduire un entretien client, qui reprend un livrable annoté,
qui entend un manager expliquer pourquoi telle formulation a été écartée. Elle ne passe pas
par le document mais par la coprésence dans l'activité. L'\textit{internalisation} ensuite~:
le savoir explicite ne devient jugement personnel qu'au prix d'une pratique répétée,
c'est-à-dire de l'apprentissage par la pratique (\textit{learning by doing}) au sens
ordinaire du terme.

C'est ici que se comprend la fonction réelle des tâches traditionnellement confiées aux
consultants débutants. Ces tâches sont bien identifiées~:
\begin{itemize}
    \item la recherche documentaire et la veille réglementaire~;
    \item la collecte, le tri et le traitement des données brutes de mission~;
    \item la première synthèse de rapports et la rédaction de livrables intermédiaires.
\end{itemize}

Perçues comme chronophages et à faible valeur ajoutée immédiate, elles remplissent en
réalité une fonction pédagogique critique. Elles obligent le débutant à se confronter à la
matière première désordonnée avant qu'elle ne soit mise en forme, et c'est cette
confrontation qui construit le jugement. Le junior qui a dû lire lui-même un texte
réglementaire pour en extraire l'exigence pertinente acquiert quelque chose que ne lui
donnera pas la lecture d'une synthèse déjà faite~: la connaissance de ce que le texte ne
dit pas, de ses ambiguïtés, des endroits où l'interprétation est ouverte. En termes du
modèle SECI, il ne se contente pas de recevoir de l'explicite~: il internalise.

L'hypothèse théorique que ce chapitre soumet à l'épreuve du terrain peut dès lors être
formulée sans ambiguïté. Un outil capable de produire instantanément la synthèse, le
premier jet du livrable ou le tableau comparatif ne supprime pas une tâche à faible valeur
ajoutée~: il supprime le support matériel d'un mode de conversion des connaissances. Le
livrable produit par la machine court-circuite l'internalisation, et le temps ainsi libéré
n'est réinvesti en socialisation que si l'organisation l'organise délibérément. Rien ne
garantit qu'elle le fasse.

\section{La fabrique de la légitimité de l'expert}

La section précédente a établi ce qu'un cabinet produit et comment il forme ceux qui
produisent. Elle laisse ouverte la question de savoir pourquoi un client accepte de payer,
puis de suivre, l'avis d'un professionnel dont il ne peut vérifier le travail. Cette
question relève de la sociologie des professions et de la théorie de la légitimité.

\subsection{La juridiction professionnelle et ses frontières}

Abbott \cite{abbott1988} propose de considérer les professions non comme des entités
isolées mais comme un système en concurrence permanente. Une profession n'existe que par la
\textit{juridiction} qu'elle parvient à revendiquer et à faire reconnaître, c'est-à-dire par
le lien exclusif qu'elle établit entre un domaine de problèmes et son propre savoir
abstrait. Ce lien s'établit par trois opérations caractéristiques~: le diagnostic, qui
classe un cas dans les catégories de la profession~; l'inférence, qui raisonne sur ce cas à
partir du savoir professionnel~; et le traitement, qui prescrit une action. La juridiction
n'est jamais acquise~: elle est disputée, devant les tribunaux, devant l'opinion et sur le
lieu de travail, par d'autres professions qui revendiquent le même territoire.

L'apport décisif de ce cadre est de désigner l'inférence comme le maillon vulnérable.
Abbott observe que les professions dont le diagnostic ou le traitement sont routinisables
perdent progressivement du terrain, tandis que la maîtrise de l'inférence — le raisonnement
sur les cas difficiles, ceux que la routine ne couvre pas — constitue le noyau défendable
d'une juridiction. Transposé au conseil, ce résultat oriente directement la troisième
question directrice. Si une technologie prend en charge une part croissante de la collecte
et de la mise en forme, la juridiction du consultant n'est pas nécessairement menacée~; elle
est déplacée vers ce qui reste non routinisable. Encore faut-il que le cabinet sache
identifier et revendiquer ce noyau, et que le client y reconnaisse une valeur. La question
n'est pas seulement de savoir ce que la machine sait faire, mais quel territoire la
profession parvient à redéfinir comme sien.

\subsection{Les trois registres de la légitimité}

Suchman \cite{suchman1995} fournit la grille permettant de décomposer ce que recouvre la
confiance accordée à un expert. La légitimité y est définie comme la perception
généralisée que les actions d'une entité sont désirables ou appropriées au regard d'un
système socialement construit de normes et de croyances. Trois registres sont distingués.
% Pagination à compléter par Arthur : tableau des registres dans \cite{suchman1995}.

La \textit{légitimité pragmatique} repose sur l'intérêt direct de l'audience~: le client
accorde du crédit parce qu'il en retire un bénéfice tangible, un problème résolu, un délai
tenu, une conformité obtenue. La \textit{légitimité cognitive} repose sur la
compréhensibilité et sur le caractère allant de soi de l'entité~: le consultant est crédible
parce que ce qu'il dit s'inscrit dans un cadre de référence partagé, parce que sa démarche
est intelligible et parce que le recours à un cabinet dans une telle situation va de soi. La
\textit{légitimité morale} repose sur un jugement normatif~: l'entité est légitime parce que
sa conduite est jugée conforme à ce qui doit être fait, indépendamment du bénéfice retiré.

Cette décomposition est opératoire pour le terrain étudié, car elle laisse prévoir que les
trois registres ne sont pas affectés de la même manière par l'automatisation. La légitimité
pragmatique peut se trouver renforcée~: un livrable plus rapide et plus complet sert
l'intérêt immédiat du client. La légitimité cognitive est ambivalente~: la production
accélérée de documents formellement irréprochables peut soutenir l'impression de maîtrise,
mais l'uniformisation des productions peut aussi éroder le caractère distinctif de
l'expertise. La légitimité morale est la plus exposée, parce qu'elle engage la question de
ce qui a été réellement fait et par qui. Un client qui découvre qu'une part substantielle de
l'analyse facturée a été produite sans intervention humaine significative ne conteste pas
d'abord la qualité du résultat~: il conteste la conformité de la conduite du cabinet à ce
qu'il croyait acheter.

\subsection{Ce sur quoi repose la confiance dans une prestation intellectuelle}

La jonction entre ces deux cadres et la section 2.1 s'opère par l'ambiguïté de l'output déjà relevée.
Si le résultat d'une mission de conseil n'est pas
directement évaluable, alors la légitimité n'est pas un supplément d'âme ajouté à la
prestation~: elle en constitue une part du produit lui-même \cite{alvesson2001}. Le client
n'achète pas seulement un diagnostic~; il achète l'assurance que ce diagnostic a été produit
par quelqu'un qui en répond.

Cette formulation permet de préciser la question centrale du mémoire. Ce qui est en jeu
lorsqu'une machine produit une part substantielle de la valeur intellectuelle n'est pas la
disparition de l'expert, mais le déplacement du fondement de son autorité. L'autorité
adossée au volume de travail fourni et à la maîtrise technique de la collecte s'affaiblit
mécaniquement. Celle qui s'adosse à la responsabilité assumée sur un jugement — au sens de
l'inférence d'Abbott et de la légitimité morale de Suchman — devient, en théorie, le
territoire défendable. Que ce déplacement soit effectivement perçu, revendiqué et valorisé
par les consultants et par leurs clients relève du terrain, et non de la théorie.

\section{L'intelligence artificielle générative comme technologie de rupture dans les
systèmes d'information}

Les deux sections précédentes ont construit la structure et la légitimité. Il reste à
caractériser la technologie qui les percute, sous un angle strictement organisationnel~: ce
qui importe ici n'est pas la manière dont ces outils fonctionnent, mais la manière dont ils
sont adoptés, ce qu'ils déplacent dans la division du travail et ce qu'ils font au jugement
professionnel.

\subsection{Les modèles d'acceptation technologique et leurs limites}

Le modèle d'acceptation de la technologie (\textit{Technology Acceptance Model}) formulé par
Davis \cite{davis1989} explique l'usage effectif d'un système d'information par deux
variables perçues~: l'utilité perçue et la facilité d'utilisation perçue. Sa longévité tient
à sa parcimonie. Ses prolongements successifs en ont progressivement corrigé le caractère
individualiste. Le TAM2 \cite{venkatesh2000} introduit l'influence sociale et les processus
instrumentaux cognitifs~: l'usage d'un pair, la norme subjective et l'image tirée de
l'adoption pèsent sur la décision d'utiliser. Le modèle unifié UTAUT \cite{venkatesh2003}
ajoute aux déterminants de l'intention les \textit{conditions facilitatrices},
c'est-à-dire l'ensemble des ressources organisationnelles rendant l'usage possible. Le TAM3
\cite{venkatesh2008} franchit un pas supplémentaire en identifiant les interventions
organisationnelles susceptibles d'agir sur l'acceptation~: soutien managérial, formation,
participation des utilisateurs.

Ces prolongements servent la première et la quatrième question directrices. Le TAM2 éclaire
la diffusion par imitation entre pairs, plausible dans des équipes projet où l'usage d'un
outil se transmet par observation plutôt que par prescription. L'UTAUT désigne le
management comme producteur des conditions facilitatrices, ce qui est exactement l'objet de
la question de gouvernance. Le TAM3 fournit le répertoire des leviers actionnables et
constitue à ce titre l'appui théorique des recommandations formulées ultérieurement.

Deux réserves doivent cependant être posées. La première tient à l'objet~: ces modèles ont
été construits sur des logiciels de bureautique dont l'usage était prescrit par
l'employeur, non sur des outils que le collaborateur peut mobiliser de sa propre
initiative, hors du cadre officiel, y compris depuis un terminal personnel. La distinction
entre adoption prescrite et appropriation spontanée, que le TAM ne thématise pas, est
centrale sur le terrain étudié. La seconde tient à la variable expliquée~: ces modèles
expliquent l'intention d'usage, non les effets de l'usage sur la compétence de
l'utilisateur. Or c'est précisément l'effet, et non l'adoption, qui est ici en question. Le
cadre de l'acceptation est donc nécessaire mais insuffisant~; il doit être complété par les
travaux empiriques portant sur les conséquences de l'usage.

\subsection{Du consultant augmenté au consultant assisté~: ce qu'établissent les travaux
empiriques}

L'expérimentation la plus proche du terrain de ce mémoire a été conduite en 2023 auprès de
consultants du Boston Consulting Group, en collaboration avec le centre de recherche du
groupe, et publiée depuis dans \textit{Organization Science}
\cite{dellacqua2026}. Son résultat central est la notion de \textit{frontière dentelée}
(\textit{jagged frontier})~: le périmètre des tâches sur lesquelles l'outil apporte un gain
n'est ni continu ni prévisible depuis l'extérieur. Sur les tâches situées à l'intérieur de
ce périmètre, les consultants équipés produisent davantage, plus vite et avec une qualité
supérieure. Sur des tâches d'apparence voisine mais situées au-delà de la frontière, l'usage
de l'outil dégrade la performance~: les consultants assistés se trompent plus souvent que
ceux qui travaillent sans assistance.

Ce résultat est le plus important du chapitre, pour deux raisons. D'abord parce qu'il
interdit de poser la question en termes de gain global~: l'effet est conditionnel à la
nature de la tâche, et la difficulté pratique tient à ce que la frontière est invisible pour
celui qui l'approche. Ensuite parce qu'il déplace le problème managérial. Si la frontière
était identifiable, il suffirait de prescrire l'usage en deçà et de l'interdire au-delà.
Comme elle ne l'est pas, la seule ressource disponible est le jugement du professionnel sur
sa propre situation — c'est-à-dire exactement la compétence dont la section 2.1 laisse
craindre l'érosion.

Deux travaux complémentaires portent sur la distribution des gains. Noy et Zhang
\cite{noy2023} établissent, sur des tâches de rédaction professionnelle, un gain de
productivité assorti d'un resserrement de l'écart entre les participants les moins
performants et les plus performants~: l'outil profite davantage aux premiers. Brynjolfsson,
Li et Raymond \cite{brynjolfsson2025} retrouvent ce schéma en conditions réelles de
déploiement en entreprise, où les gains se concentrent sur les travailleurs les moins
expérimentés. La convergence entre un protocole expérimental et des données d'entreprise
donne à ce résultat une robustesse peu commune dans la littérature récente.

Sa lecture, en revanche, est ambivalente, et cette ambivalence structure la deuxième
question directrice. Une lecture optimiste y voit une accélération de la montée en
compétences~: le junior atteint plus tôt un niveau de production acceptable. Une lecture
pessimiste observe que le resserrement porte sur la production, non sur la compétence~: si
le débutant atteint le résultat du senior sans avoir parcouru le chemin qui produit le
jugement, l'écart de performance mesurée masque un écart de compétence inchangé, voire
creusé. La théorie ne permet pas de trancher entre ces deux lectures~; le terrain, oui.
C'est précisément ce que le dispositif d'entretiens du chapitre suivant cherche à établir en
confrontant le ressenti des juniors sur leur propre apprentissage à la perception qu'en ont
leurs encadrants.

Enfin, les travaux d'Eloundou, Manning, Mishkin et Rock \cite{eloundou2024} sur l'exposition
des métiers aux modèles de langage situent ce mouvement à l'échelle macro-économique et
confirment que les professions intellectuelles à forte composante rédactionnelle et
analytique figurent parmi les plus exposées. Cette référence sert au cadrage général~; elle
ne dit rien de ce qui se joue à l'échelle d'une équipe.

\subsection{Opacité, contrôle et atrophie de la réflexivité critique}

Le dernier ensemble de travaux porte sur ce que l'usage d'une technologie opaque fait au
jugement professionnel. Anthony \cite{anthony2021} conduit une étude ethnographique
d'analystes juniors confrontés à une technologie analytique dont le fonctionnement leur
échappe. Le résultat est nuancé et d'autant plus utile~: le clivage ne passe pas entre ceux
qui utilisent l'outil et ceux qui s'en abstiennent, mais entre deux manières de l'utiliser.
Les analystes qui interrogent la production de l'outil, qui cherchent à comprendre d'où
vient un résultat et à en éprouver la cohérence, développent une expertise plus solide. Ceux
qui l'utilisent sans l'interroger — parce que le résultat paraît plausible, parce que le
délai presse, parce que l'outil fait autorité — développent une expertise plus
superficielle. Le mécanisme est identifié~: c'est l'acceptation non critique du résultat, et
non l'usage lui-même, qui appauvrit l'apprentissage.

Ce résultat est le plus directement transposable au cas étudié. Il désigne une variable
managériale actionnable, là où le constat d'une atrophie générale des compétences serait
resté sans prise. Il fournit également un critère d'analyse pour le chapitre~4~: la question
posée aux répondants n'est pas de savoir s'ils utilisent ces outils, mais ce qu'ils font du
résultat produit, et à quelles conditions ils le remettent en cause.

Lebovitz, Levina et Lifshitz-Assaf \cite{lebovitz2021} apportent un argument de fond
souvent absent des discussions sur la fiabilité de ces outils. Les « vérités terrain » sur
lesquelles une technologie d'intelligence artificielle est entraînée et évaluée sont
elles-mêmes construites et contestables~; l'évaluation d'un outil suppose donc un étalon qui
n'est pas toujours disponible. Appliquée au conseil, la portée de cet argument est
considérable~: le problème posé par les productions erronées n'est pas seulement que l'outil
se trompe, c'est que le professionnel ne dispose pas toujours d'une référence indiscutable
contre laquelle vérifier. Le contrôle qualité ne peut donc être conçu comme une simple
opération de vérification factuelle.

Les mêmes auteurs \cite{lebovitz2022} examinent comment des professionnels arbitrent
concrètement face à cette opacité lorsqu'ils doivent porter un jugement engageant leur
responsabilité. Le parallèle avec le consultant qui signe un livrable remis au client est
direct. Cette étude constitue l'appui principal de la quatrième question directrice et la
matrice des protocoles de validation qui seront proposés en recommandation.

Kellogg, Valentine et Christin \cite{kellogg2020} complètent le dispositif en analysant le
contrôle algorithmique comme un nouveau terrain de lutte dans l'organisation. Les auteurs
recensent les leviers par lesquels une technologie oriente, évalue et discipline le travail,
ainsi que les formes de résistance que ces leviers suscitent. Cette grille est
indispensable à la formulation des recommandations de gouvernance, car elle rappelle
qu'encadrer un usage n'est jamais un acte neutre~: toute règle produit des effets
d'appropriation et de contournement. Elle éclaire directement la tension centrale de la
quatrième question directrice, entre la nécessité de maîtriser un risque opérationnel et le
maintien de l'autonomie professionnelle sans laquelle une bureaucratie professionnelle ne
fonctionne pas.

\subsection{Synthèse~: quatre questionnements, quatre cadres d'analyse}

L'ensemble des cadres mobilisés peut être rapporté aux quatre questionnements directeurs du
mémoire, chacun trouvant ici les instruments de son traitement empirique.

\begin{table}[h]
\centering
\small
\begin{tabular}{|p{0.24\linewidth}|p{0.33\linewidth}|p{0.33\linewidth}|}
\hline
\textbf{Questionnement} & \textbf{Cadre théorique mobilisé} & \textbf{Appui empirique} \\
\hline
1. Appropriation réelle des outils
& TAM et prolongements~: influence sociale, conditions facilitatrices
(\textsc{Davis}, 1989~; \textsc{Venkatesh} et \textsc{Davis}, 2000~;
\textsc{Venkatesh} \textit{et al.}, 2003)
& Distribution non uniforme des usages et des gains
(\textsc{Dell'Acqua} \textit{et al.}, 2026) \\
\hline
2. Apprentissage et montée en compétences des juniors
& Bureaucratie professionnelle et standardisation des qualifications
(\textsc{Mintzberg}, 1982)~; conversion des connaissances et savoir tacite
(\textsc{Polanyi}, 1966~; \textsc{Nonaka} et \textsc{Takeuchi}, 1997)
& Resserrement des écarts de performance (\textsc{Noy} et \textsc{Zhang}, 2023~;
\textsc{Brynjolfsson} \textit{et al.}, 2025)~; expertise superficielle en cas d'usage non
critique (\textsc{Anthony}, 2021) \\
\hline
3. Légitimité de l'expert face au client
& Juridiction professionnelle et inférence (\textsc{Abbott}, 1988)~; registres
pragmatique, cognitif et moral (\textsc{Suchman}, 1995)~; ambiguïté de l'\textit{output}
(\textsc{Alvesson}, 2001) & Déplacement de la valeur vers les tâches non routinisables
(\textsc{Dell'Acqua} \textit{et al.}, 2026) \\
\hline
4. Gouvernance et maîtrise des risques
& Contrôle algorithmique et résistances (\textsc{Kellogg} \textit{et al.}, 2020)~;
interventions organisationnelles sur l'acceptation (\textsc{Venkatesh} et \textsc{Bala},
2008) & Jugement professionnel en situation d'opacité (\textsc{Lebovitz} \textit{et al.},
2021 et 2022) \\
\hline
\end{tabular}
\caption{Correspondance entre les questionnements directeurs et les cadres d'analyse
mobilisés}
\label{tab:cadres-questionnements}
\end{table}

Le tableau~\ref{tab:cadres-questionnements} rapporte l'ensemble des cadres mobilisés aux quatre questionnements directeurs. Les cadres retenus ne sont pas juxtaposés~: ils se
répondent. La bureaucratie professionnelle explique pourquoi une perturbation de
l'apprentissage est une question structurelle et non une question de formation~; la théorie
de la légitimité explique pourquoi cette perturbation finit par atteindre la relation
client~; les travaux empiriques sur l'usage établissent que l'effet dépend de la manière
d'utiliser plus que de l'usage lui-même, ce qui redonne au management une prise. La chaîne
de raisonnement soumise au terrain peut être énoncée en une phrase~: l'automatisation des
tâches d'analyse déplace le lieu où se forme le jugement professionnel, et la question
managériale est de savoir si l'organisation reconstitue ailleurs ce qu'elle supprime là.

Le chapitre suivant expose la démarche retenue pour soumettre cette proposition à
l'épreuve des faits.

\cleardoublepage{}

% ==============================================================================
% PARTIE 2 : CADRE EMPIRIQUE
% ==============================================================================
\part{Cadre Empirique -- Méthodologie et Réalité du terrain}

\chapter{Présentation du terrain et environnement macro-économique}
%% ==============================================================================
% CHAPITRE 3 -- DÉMARCHE MÉTHODOLOGIQUE QUALITATIVE
% Mémoire MAE -- Arthur SAUNIER -- Wavestone Luxembourg
% Le \chapter{} est déclaré dans le fichier maître ; ce fichier démarre en \section.
% ==============================================================================

\section{Justification du design de recherche}

La question centrale de ce mémoire interroge la manière dont l'automatisation des
tâches d'analyse redéfinit la légitimité et le développement des compétences des
consultants. Sa formulation en « comment » n'appelle ni la mesure d'un écart ni la
vérification d'une relation entre variables~: elle appelle la compréhension d'un
processus en cours, saisi dans les pratiques et dans les représentations de ceux qui
le vivent. C'est cette nature de la question qui commande le choix méthodologique, et
non l'inverse.

Une démarche quantitative supposerait que les dimensions à mesurer soient déjà
stabilisées. Or elles ne le sont pas. Construire un questionnaire reviendrait à
décider par avance ce qu'est un usage, ce qu'est une perte de compétence ou ce qu'est
une atteinte à la légitimité, puis à demander aux répondants de se positionner sur des
catégories dont rien ne garantit qu'elles recouvrent leur expérience. L'objet est
émergent~; les outils concernés se sont diffusés dans les pratiques plus vite que le
vocabulaire permettant de les décrire. Une enquête par questionnaire produirait donc
une mesure précise de catégories mal fondées.

Deux obstacles pratiques renforcent ce choix. Le premier tient à la taille de la
population~: l'effectif de la practice concernée ne permet aucune exploitation
statistique digne de ce nom, la puissance d'un test étant hors d'atteinte sur quelques
dizaines d'individus. Le second tient à la sensibilité du sujet. Les usages
informels, le recours à des outils non validés par le cabinet ou la validation d'un
contenu que l'on n'aurait pas su produire soi-même ne se déclarent pas dans une case à
cocher. Ces pratiques ne s'obtiennent que par le récit, dans un cadre où le répondant
comprend que la description d'un contournement n'est pas un aveu.

La posture épistémologique retenue est interprétativiste. Elle admet que la réalité
étudiée n'est pas donnée indépendamment des acteurs qui la construisent, et que le
matériau accessible au chercheur est le sens que ces acteurs attribuent à leurs propres
pratiques \cite{gavardperret2018}. Le discours recueilli n'est donc pas traité comme
un reflet transparent de l'activité réelle, mais comme une reconstruction située, dont
les silences, les hésitations et les contradictions internes font partie du matériau
plutôt qu'ils n'en altèrent la qualité.

Le raisonnement suivi est abductif. Le cadrage théorique du chapitre précédent fournit
les catégories de départ~: modes de conversion des connaissances, registres de
légitimité, mécanismes de coordination, déterminants de l'acceptation. Ces catégories
orientent la construction de l'instrument de collecte sans prétendre épuiser ce que le
terrain donnera à voir. L'allée et venue entre les concepts mobilisés et le matériau
recueilli constitue le mouvement analytique propre à ce travail.

L'entretien semi-directif individuel a été préféré aux autres techniques qualitatives
disponibles. Le groupe de discussion aurait exposé les propos les plus sensibles à la
censure des pairs et à la présence hiérarchique. L'observation directe aurait mal
saisi des usages qui se déroulent pour l'essentiel dans un rapport solitaire à un
écran. L'entretien semi-directif, en revanche, articule une liberté de formulation
suffisante pour laisser émerger l'inattendu et une structure commune sans laquelle la
comparaison entre profils serait impossible \cite{blanchet2007}.

Il importe enfin d'énoncer ce que ce design autorise et ce qu'il interdit. Il autorise
l'identification de mécanismes, la mise au jour de tensions entre niveaux
hiérarchiques et la formulation d'hypothèses transférables à des organisations
comparables. Il interdit toute généralisation statistique~: aucun résultat ne pourra
être présenté comme valant pour la profession du conseil, ni même pour l'ensemble du
cabinet. Il interdit également toute conclusion causale et toute mesure de
performance. Si un répondant décrit un gain de temps, ce gain de temps est un fait de
discours et non un fait mesuré~; le dispositif établit ce qui est perçu, vécu et
raconté, non ce qui est produit. Cette frontière est rappelée à chaque étape de la
restitution.

\section{Dispositif empirique et échantillonnage}

L'échantillon visé compte six à huit collaborateurs de Wavestone Luxembourg, tous
rattachés à la practice Cybersécurité. Il ne procède d'aucun tirage aléatoire mais
d'un échantillonnage raisonné, construit par contraste de profils. La variable
discriminante retenue est la position hiérarchique, parce qu'elle détermine le rapport
même à l'objet étudié~: le consultant junior éprouve sa propre montée en compétences,
le manager l'observe chez ceux qu'il encadre, le directeur en arbitre les conséquences
sur le modèle de production du cabinet. Un même phénomène est ainsi saisi depuis trois
positions d'énonciation distinctes.

\begin{table}[h]
\centering
\small
\begin{tabular}{|p{0.13\linewidth}|p{0.09\linewidth}|p{0.32\linewidth}|p{0.32\linewidth}|}
\hline
\textbf{Profil} & \textbf{Codes} & \textbf{Position dans le dispositif} & \textbf{Apport attendu} \\
\hline
Consultants juniors & J1--J3 & Producteurs directs des livrables d'analyse & Usage vécu, sentiment sur l'apprentissage, rapport aux règles \\
\hline
Seniors et managers & M1--M3 & Encadrants de production et responsables du contrôle qualité & Observation de l'évolution des juniors, rôle en revue, relation client directe \\
\hline
Directeurs & D1--D2 & Responsables du modèle économique et de la politique d'outillage & Arbitrage risque-productivité, proposition de valeur, staffing et formation \\
\hline
\end{tabular}
\caption{Composition du panel et logique de contraste}
\label{tab:panel-chap3}
\end{table}

Le tableau~\ref{tab:panel-chap3} explicite la logique de construction du panel. Trois
juniors, trois seniors ou managers et un à deux directeurs permettent de constituer
deux groupes comparables en effectif, condition d'une lecture croisée, tout en
réservant une position de surplomb. Le tiraillement attendu entre le groupe des
juniors et celui des managers est une propriété voulue du design et non un résultat~:
il conditionne la possibilité même d'observer un écart, mais ne préjuge en rien de son
existence ni de son sens. Des critères secondaires --- ancienneté dans le cabinet,
diversité des types de missions --- guident le choix des individus à l'intérieur de
chaque strate, afin d'éviter que le panel ne se réduise à un même segment d'activité.

Le guide d'entretien, reproduit intégralement en annexe~A, est construit de manière
déductive à partir des quatre questions directrices. Chaque thème constitue la
traduction conversationnelle d'une de ces questions. Un thème de cadrage préalable,
consacré à la trajectoire du répondant et à la nature de ses missions, ouvre
l'entretien~: il installe la parole, recueille les variables de contrôle nécessaires à
la lecture croisée et permet à l'enquêteur de repérer le vocabulaire propre au
répondant, qu'il réemploiera ensuite. Le thème~1 porte sur l'usage réel des outils
dans le travail d'analyse, le thème~2 sur l'apprentissage du métier et la construction
de l'expertise, le thème~3 sur la relation au client et la valeur de l'expertise, le
thème~4 sur l'encadrement, le contrôle qualité et la gestion des risques. Un thème de
clôture, projectif et ouvert, invite le répondant à formuler une décision qu'il
prendrait à la place de la direction et à signaler ce que l'entretien n'a pas abordé.

Deux partis pris commandent la formulation des questions \cite{blanchet2007}. Chaque
thème s'ouvre par un récit ou un incident critique plutôt que par une question
d'opinion, le récit d'une situation vécue produisant un matériau vérifiable là où
l'opinion appelle la réponse de façade. L'enquêteur s'interdit par ailleurs de
prononcer le vocabulaire de la recherche --- légitimité, montée en compétences,
hallucination, gouvernance --- avant que le répondant ne l'ait lui-même mobilisé, tout
emploi prématuré induisant la réponse et détruisant la valeur probante du verbatim. La
structure du guide demeure identique pour tous les répondants, condition de la
comparabilité~; seules les formulations d'ancrage sont adaptées au profil, afin de
solliciter la position d'énonciation propre à chacun.

L'entretien est calibré sur trente-cinq à quarante-cinq minutes, durée compatible avec
les contraintes de disponibilité de consultants en mission. Les entretiens se tiennent
en visioconférence ou en présentiel, mais toujours au moyen d'une réunion organisée
dans l'environnement collaboratif du cabinet~: l'enregistrement et la transcription
automatique sont assurés par cet outil, ce qui garantit que l'audio ne quitte jamais
le système d'information de l'organisation. Une seule langue est retenue par entretien,
l'alternance linguistique dégradant fortement la qualité de la transcription
automatique.

\section{Modalités de collecte et d'analyse}

Chaque entretien s'ouvre par la lecture d'un préambule reproduit en annexe~A. Ce
préambule précise l'objet du travail, rappelle qu'il ne s'agit ni d'un audit ni d'une
évaluation, énonce les modalités d'anonymisation et sollicite explicitement
l'autorisation d'enregistrer. En cas de refus, l'entretien se poursuit sous prise de
notes manuscrite et la restitution des propos est alors signalée comme reconstituée.
Les répondants sont désignés par un code de profil suivi d'un rang~; les noms de
personnes, de clients et de missions sont remplacés par des désignations génériques
entre crochets dès la retranscription. Une fiche de traçabilité est renseignée
immédiatement après chaque entretien, avant toute retranscription, afin de conserver
les éléments contextuels que l'enregistrement ne capte pas~: climat de l'échange,
réticences perçues, thèmes écourtés, hypothèses émergentes à tester lors des entretiens
suivants.

La transcription automatique produite par l'outil de visioconférence fait l'objet d'une
reprise manuelle intégrale. Cette relecture n'est pas une simple correction
orthographique~: elle rétablit les hésitations, les reprises et les ruptures de
construction, qui sont autant d'indices analytiques. Le corpus final est constitué des
transcriptions ainsi corrigées, à l'exclusion de toute reformulation.

Le traitement repose sur une analyse thématique de contenu \cite{bardin2001}. Le
matériau est découpé, catégorisé, puis interprété selon les trois temps classiques de
la préanalyse, de l'exploitation et de l'inférence. L'unité de codage retenue est le
segment de sens, c'est-à-dire l'extrait de discours porteur d'une idée complète,
indépendamment de sa longueur grammaticale. Un même segment peut recevoir plusieurs
codes lorsqu'il articule plusieurs dimensions, un consultant décrivant par exemple dans
la même phrase un gain de productivité et une inquiétude sur sa propre compétence.

La grille de codage, reproduite en annexe~B, est mixte. Son ossature est déductive~:
quatre axes correspondent terme à terme aux quatre questions directrices et à leurs
ancrages théoriques respectifs. L'axe~A traite de l'appropriation des outils dans le
travail d'analyse, en distinguant les tâches déléguées, les tâches sanctuarisées, les
canaux d'apprentissage de l'outil et les déterminants d'usage. L'axe~B traite de
l'apprentissage, des compétences et de la transmission, en distinguant la valeur
formatrice attribuée aux tâches, les effets déclarés sur le raisonnement et les
transformations de la relation entre junior et senior. L'axe~C traite de la légitimité
et de la relation client, en distinguant les registres de confiance invoqués, le
traitement de la transparence vis-à-vis du client et les déplacements de la
proposition de valeur. L'axe~D traite de la gouvernance, du contrôle qualité et des
risques, en distinguant le statut des règles, les dispositifs de vérification, les
incidents rencontrés, l'attribution des responsabilités et la perception du cadre
lui-même, jugé tantôt bridant, tantôt insuffisamment explicite. Un cinquième axe,
transversal, est partiellement inductif~: il accueille les écarts entre discours et
pratique, les divergences entre niveaux hiérarchiques, les projections sur l'avenir du
métier, les recommandations spontanées et une catégorie ouverte destinée aux thèmes
non anticipés. Ces derniers sont intitulés et définis en cours d'analyse.

Le codage est conduit au fil de l'eau, chaque entretien étant traité avant la
conduite du suivant. Cette contrainte n'est pas d'ordre pratique~: elle est la
condition d'un suivi honnête de la saturation théorique \cite{glaser1967}. La grille
est stabilisée après le troisième entretien, puis appliquée rétroactivement à
l'ensemble du corpus afin que tous les entretiens soient codés selon le même
référentiel. La saturation est considérée comme atteinte lorsqu'un entretien
supplémentaire n'ajoute aucun code nouveau et ne modifie pas la structure de la
grille~; le rang de l'entretien à partir duquel aucun code inédit n'apparaît est
mentionné dans la restitution. Les contradictions internes à un même entretien ne sont
pas arbitrées mais codées comme telles~: elles constituent un résultat.

Le corpus codé alimente une matrice de restitution croisant les codes en lignes et les
répondants en colonnes, complétée d'une colonne de fréquence et d'un verbatim
représentatif \cite{miles2003}. La lecture en colonnes restitue la cohérence interne
du discours de chaque répondant~; la lecture en lignes fait apparaître les
convergences du panel et ses lignes de fracture, une colonne d'écart signalant les
codes dont la fréquence diffère sensiblement entre le groupe des juniors et celui des
managers. Cette matrice constitue le support direct du chapitre suivant~: chaque
verbatim y cité est référencé par le code du répondant et le code thématique, de
manière à ce que le chemin analytique entre le corpus brut et l'interprétation proposée
demeure reconstituable par le lecteur.

\section{Limites et biais du dispositif}

La taille de l'échantillon constitue la première limite. Six à huit entretiens
autorisent la profondeur, non la représentativité. Rien ne garantit d'ailleurs que la
saturation soit effectivement atteinte au terme du panel prévu~; si des codes inédits
continuent d'apparaître au dernier entretien, ce fait sera signalé comme tel plutôt que
dissimulé, et les conclusions correspondantes présentées comme provisoires.

Le caractère mono-site constitue la deuxième limite. L'étude porte sur une practice
d'une seule entité, sur un marché dont la structure --- place financière concentrée,
pression réglementaire spécifique --- n'est pas transposable telle quelle. Les résultats
sont indexés à ce contexte. Leur portée relève de la transférabilité analytique vers
des organisations comparables, non de la généralisation \cite{lincoln1985}.

La désirabilité sociale constitue la troisième limite, et elle est ici plus vive
qu'ailleurs. Le sujet touche à des pratiques qu'un collaborateur n'a pas intérêt à
exposer~: le recours à des outils hors du cadre officiel, le contournement d'une règle,
l'aveu d'une dépendance ou la validation d'un contenu que l'on n'aurait pas su produire.
Plusieurs dispositions visent à contenir ce biais~: le préambule dissocie explicitement
l'entretien de toute démarche d'évaluation, l'anonymisation est décrite avant
l'enregistrement, les questions procèdent par récit d'incident plutôt que par
sollicitation d'opinion, et l'écart entre position déclarée et pratique décrite est
codé plutôt que résolu. Ces précautions atténuent le biais sans le supprimer.

La position de l'enquêteur constitue la limite la plus lourde et appelle un examen
séparé. Celui-ci est lui-même consultant du cabinet étudié~: il occupe simultanément la
place de l'observateur et celle du participant \cite{girin1989}. Ce double statut n'est
pas un accident
du dispositif, il en est une condition d'existence, et il produit des effets dans les
deux sens.

Les avantages sont réels. L'accès au terrain est immédiat, là où un chercheur externe
se heurterait aux règles de confidentialité qui régissent l'activité de conseil. La
familiarité avec le vocabulaire professionnel, avec le déroulement d'une mission et
avec la chaîne de production d'un livrable permet de comprendre l'implicite, de relancer
au bon endroit et de repérer une réponse de façade. La confiance préexistante entre
collègues rend possible l'évocation de pratiques qu'un inconnu n'obtiendrait pas.

Les risques sont symétriques. Le répondant s'adresse à un collègue et ajuste son
propos en conséquence, particulièrement lorsque l'écart hiérarchique joue en sens
inverse de la situation d'entretien, un directeur restant en position d'autorité face à
un enquêteur qui lui est subordonné. L'enquêteur, de son côté, aborde le terrain avec
des prénotions constituées par sa propre expérience et risque de sélectionner ce qui
les confirme. La connivence professionnelle produit enfin un effet insidieux~: la
formule « tu vois ce que je veux dire » clôt l'explication au moment précis où elle
devient nécessaire, et l'enquêteur qui voit effectivement ce que le répondant veut dire
perd le verbatim qui l'aurait établi.

Quatre mesures visent à contenir ces risques. Le statut de l'enquêteur est explicité
en préambule, de sorte que l'ambiguïté de la situation soit posée plutôt que subie.
L'explicitation de l'implicite est systématiquement demandée, y compris --- et surtout ---
lorsque l'enquêteur a compris~: toute allusion est reprise par une demande de récit
factuel. Une neutralité axiologique stricte est tenue, l'enquêteur s'abstenant de
défendre comme de critiquer l'usage de ces outils, toute marque d'approbation orientant
durablement la suite de l'entretien. La frontière entre les deux statuts est enfin
maintenue dans le traitement~: seul le matériau enregistré est codé, la connaissance de
terrain servant à interpréter les propos mais jamais à établir un fait que personne
n'aurait énoncé.

Ces limites ne sont pas des défauts à corriger mais des propriétés du dispositif. Les
énoncer avant la présentation des résultats est la condition d'une lecture avertie de
ce qui suit.

\chapter{Démarche méthodologique qualitative}
%% ==============================================================================
% CHAPITRE 4 -- RESTITUTION DES RÉSULTATS, CONFRONTATION ET RECOMMANDATIONS
% Mémoire MAE -- Arthur SAUNIER -- Wavestone Luxembourg
% Le \chapter{} est déclaré dans le fichier maître ; ce fichier démarre en \section.
% ==============================================================================

\section{Restitution du matériau}

La restitution suit les axes de la grille reproduite en annexe~B et non l'ordre des
répondants. La matrice croisant les codes et les répondants \cite{miles2003} ne produit de
résultat que lue en lignes, sur ce qui converge, ce qui sépare et ce qui n'apparaît qu'une
fois. Un constat est dit \textbf{saturé} lorsque les cinq récits le documentent,
\textbf{clivant} lorsque le matériau se scinde entre deux postures, \textbf{isolé} lorsqu'un
seul répondant le rapporte. Il est alors conservé, mais sans portée générale. Le mot
«~saturé~» vaut dans son acception faible~: sur cinq entretiens codés après la clôture de la
collecte, la saturation théorique \cite{glaser1967} n'est pas revendiquée.

Le suivi du nombre de codes nouveaux apportés par chaque entretien, dans l'ordre de
passation, précise cette réserve. Le premier entretien fait apparaître vingt-sept codes,
le deuxième onze, le troisième quatre, le quatrième deux et le cinquième quatre~;
quarante-huit des cinquante-trois codes de la grille sont mobilisés au total. La
décroissance est nette et la structure de la grille se stabilise dès le troisième
entretien, les cinq axes étant tous mobilisés dès le premier récit. Aucun entretien
n'atteint toutefois zéro, et le dernier apporte encore quatre codes inédits, dont deux relèvent de l'axe~C~: la saturation n'est pas atteinte, et la moitié de ce qu'il
apporte encore relève de l'axe que la section~4.4 déclare le plus
fragile.

Le tableau~\ref{tab:statut-resultats} récapitule le statut des principaux résultats au regard du panel ; les sous-sections qui suivent les développent axe par axe.

\begin{table}[h]
\centering
\small
\begin{tabular}{|p{0.46\linewidth}|p{0.15\linewidth}|p{0.25\linewidth}|}
\hline
\textbf{Résultat} & \textbf{Codes} & \textbf{Statut} \\
\hline
Zone sanctuarisée, et convergence de sa frontière & A2.1, A2.2 & Saturé \\
\hline
Seuil au-delà duquel l'outil coûte plus qu'il ne rapporte & A4.1, A4.3 & Saturé \\
\hline
Temps gagné réinvesti dans la vérification & A4.1, B2.3 & Saturé, une exception \\
\hline
Apprentissage de l'outil autonome ou entre pairs & A3.1, A3.2 & Saturé, une exception \\
\hline
Déclaration ou non de l'usage au client & C2.2, C2.1 & Saturé, une exception \\
\hline
Valeur déplacée vers le jugement et la responsabilité & C3.1, C3.2 & Saturé \\
\hline
Absence de contrôle dédié aux contenus générés & D2.3 & Saturé \\
\hline
Cadre normatif jugé insuffisant & D1.3, D5.2 & Saturé \\
\hline
Effet perçu sur l'apprentissage du junior & B2.2, T1.2 & Clivant \\
\hline
Adaptation de la pratique de revue & D2.2, T3.2 & Clivant \\
\hline
Interpellation du cabinet par un client & C2.3 & Isolé \\
\hline
Transparence assumée vis-à-vis du client & C2.1 & Isolé \\
\hline
\end{tabular}
\caption{Statut des principaux résultats au regard du panel}
\label{tab:statut-resultats}
\end{table}

\subsection{Axe A --- L'appropriation des outils}

Les tâches déléguées sont les mêmes d'un bout à l'autre du panel et relèvent toutes du
traitement d'un matériau déjà constitué~: reformulation, synthèse, structuration d'un plan,
traduction, premier jet sur un sujet peu maîtrisé. Aucun ne décrit une délégation de
l'analyse elle-même~; la formule d'E3 résume la position commune, «~assistant, pas
rédacteur~».

La sanctuarisation est saturée, et sa frontière l'est aussi. La zone que chacun ne délègue
pas se définit par la sensibilité des données et par la nécessité d'un contexte que l'outil
n'a pas connu, non par la difficulté de la tâche. Ce sont les données chiffrées
confidentielles pour J1 («~dealbreaker~»), l'analyse conduite après un exercice de crise
pour J2, au motif noté que «~l'IA n'était pas dans la salle~». Pour E3, ce sont
l'architecture de sécurité, les vulnérabilités et les données classifiées réservées à la
«~main humaine~»~; pour E2, les sujets spécifiques et sensibles~; pour E1, les choix
engageant une recommandation ou une orientation. Les missions diffèrent~; la frontière,
elle, ne varie pas.

Le seuil de spécificité est le deuxième constat saturé, et il est chiffré par tous. J1
compte vingt minutes de reformulations de consignes sans résultat exploitable puis quinze
minutes de rédaction directe, sur un sujet où l'outil renvoyait à des textes inexistants.
J2 compte dix minutes puis cinq. Pour E3, quarante-cinq minutes de reprise sur une
reformulation où deux concepts clés avaient été intervertis, contre vingt minutes en
production directe. Pour E1, une nuit à reprendre un livrable client intégralement produit
avec assistance. E2 énonce le même seuil par son complément, estimant à «~80~\% du travail~» la part de gain net
et désignant comme compétence propre le fait de «~savoir où pertinent, où pas~».

Le troisième constat saturé contredit l'intuition ordinaire du gain de productivité~: le
temps gagné ne sort pas de la boucle de production, il se déplace de la rédaction vers la
vérification. Le déplacement est chiffré~: deux heures en trois relectures pour J1, deux
heures en recoupement croisé des sources pour E3, une heure en contrôle de cohérence pour J2,
E1 énonçant que le gain devrait y être affecté. L'exception vient du répondant le plus
utilisateur~: E2 réinvestit ce temps dans la relation client, le c\oe{}ur du métier étant
selon ses notes de comprendre et non de produire.

L'apprentissage de l'outil se fait enfin par essai personnel ou entre pairs~: essais
personnels pour J1, apprentissage sur le tas pour J2, autodidaxie pour E2, carnet de
consignes constitué par tâtonnement pour E3. Un dispositif interne existe pourtant, et il
est obligatoire~: la formation d'environ une heure décrite au chapitre~1 conditionne
l'ouverture de l'accès à l'outil homologué. Deux répondants l'évoquent~: E1, sous la forme
d'une session de sensibilisation qu'il n'a pas approfondie, et J1, sous celle d'une session
rappelant les règles d'usage de l'outil homologué et d'un webinaire qu'une mission l'a
empêché de suivre. Aucun ne la désigne comme le canal par lequel il a appris à se servir de
l'outil. Le corpus confirme ainsi par son silence ce que le chapitre~1 établit par la
description du dispositif~: l'encadrement porte sur l'entrée dans l'usage, non sur l'usage
lui-même. Une exception doit cependant être signalée, et elle vient de l'extérieur du
cabinet. E3 suit depuis plusieurs mois une certification externe, engagée de sa propre
initiative et motivée par la demande des clients. Il note qu'elle vaut autant comme document
à valoriser que comme apprentissage, et il y revient plus tard dans l'entretien, en évoquant
sa formation continue.
L'observation a une portée qui excède l'axe~: la compétence sur ces outils commence à se
certifier hors du cabinet, ce qui rejoint le déplacement de la valeur discuté en 4.2.

\subsection{Axe B --- Apprentissage et transmission}

Une règle traverse les trois encadrants, énoncée depuis trois anciennetés différentes~:
l'exécution manuelle initiale n'est pas négociable, l'accélération vient après. E1 estime que
les juniors doivent avoir fait les tâches ingrates une ou deux fois à la main, faute de quoi
ils ne sauront pas juger ce que l'outil produit. E2 décrit la même séquence et ajoute que
déléguer sans avoir jamais produit interdit d'évaluer. E3 déplace le critère sur la
définition du métier, un bon consultant étant celui qui sait ce qu'il met dans ses livrables
et non celui qui sait formuler une consigne.

L'effet sur le raisonnement du junior est en revanche clivant. J1 documente sur lui-même ce
que la grille prévoyait d'observer. Il se dit «~plus relecteur que rédacteur~» et qualifie
d'«~illusion de compétence~» la maîtrise apparente que donne un support produit vite. Il
formule son inquiétude sous forme d'échéance, se demandant s'il saura encore, dans six mois,
«~faire structure seul~». J2, à ancienneté comparable, ne partage pas ce diagnostic. La
différence tient vraisemblablement à la nature des tâches, créatives et non factuelles chez
lui. 
Côté encadrant, le constat est convergent~: E1 décrit des «~livrables très propres très vite~»
suivis d'une interrogation sur la profondeur, et des juniors «~très bons forme, mais mal
justifier fond~»~; E3 parle d'une illusion de maîtrise sans avoir creusé.

Le résultat le plus net de cet axe est négatif. L'appauvrissement de l'interaction entre
junior et encadrant, inscrit en B3.2 à titre d'hypothèse, n'est documenté par personne, et
les deux positions décrivent le mouvement inverse. J1 observe que les commentaires de
l'encadrant portent désormais sur la solidité du raisonnement plutôt que sur l'alignement des
titres, et relève lui-même le caractère paradoxal de cette amélioration de la relecture
senior. E1 décrit la relecture comme transmission du jugement plutôt que correction
technique, affirme que cette transmission devient plus importante avec l'outil, et range la
formation des juniors à la vérification critique parmi les composantes importantes de son
rôle. Ce résultat est repris en 4.2.

\subsection{Axe C --- Légitimité et relation client}

La question de la déclaration de l'usage au client est documentée par les cinq récits, et
quatre y répondent par la non-déclaration, avec la même
justification~: le client achète un résultat, non un processus. J1 parle de sujet tabou~; J2
juge la mention sans apport. E3 la juge non pas dissimulatrice mais contre-productive, le
client payant une expertise et non un outil. E1, interrogé par des clients, répond que
l'analyse finale et les recommandations reviennent aux équipes, ce qui est exact. Mais il
indique qu'il ne dirait jamais qu'un support a été produit avec l'outil. E2 est isolé, et cet
isolement est cohérent avec sa fonction~: il aborde le sujet de sa propre initiative en
rendez-vous commercial et en fait un argument de positionnement.

L'interpellation par le client est le point le plus mince du corpus. Un seul répondant, E1,
rapporte une situation où un client a soulevé la question, et une seule conséquence
commerciale d'une défaillance est documentée~: des erreurs factuelles non détectées, une
méfiance non exprimée directement, une mission suivante sur laquelle un autre consultant a
été demandé. La question projective sur l'horizon de dix ans produit en revanche une
convergence complète, les cinq plaçant la valeur future du cabinet du côté du jugement et de
la responsabilité assumée. Hormis le récit d'E1, tout cet axe relève du déclaratif projectif.

\subsection{Axe D --- Gouvernance, contrôle et risques}

Ce que les répondants connaissent est un cadre de confidentialité et une liste d'outils
homologués, non une doctrine d'usage. Le chapitre~1 rappelle qu'il n'existe pas de politique
écrite en la matière. Tous le jugent insuffisant, pour deux raisons opposées. Pour les uns,
la règle est floue sur son périmètre~: J1 la décrit comme un flou total s'agissant des
outils non homologués et qualifie ce flou de «~dangereux~», E1 juge les règles
«~un peu légères~». Pour les autres, elle est claire mais muette sur l'essentiel. E2, qui la
trouve «~trop timide~», observe qu'elle autorise et interdit des outils sans rien dire de la
manière de vérifier ce qu'ils produisent. E3 la trouve générale au point d'appliquer une
règle plus stricte. Un cadre
existe donc~: il délimite des outils, il ne délimite pas des pratiques.

L'absence de contrôle dédié suit. Le circuit de revue hiérarchique fonctionne et personne ne
le conteste, mais aucun des cinq ne décrit d'étape propre aux contenus produits avec
assistance. J1 est le plus explicite~: personne ne demande d'où vient un contenu, le contrôle
porte sur le fond et non sur le processus, et les encadrants «~soupçonnent mais ne posent
pas~». E3 confirme qu'aucune procédure spécifique n'existe et rappelle que la charge repose
sur celui qui signe.

Les incidents, eux, existent~: un livrable client entier comportant chiffres faux et
références inventées, repris en une nuit, pour E1~; plusieurs occurrences pour E3~; un
quasi-incident pour J1, sur un document individuel où des projets inexistants avaient été
introduits. Leur traitement constitue le résultat de gouvernance le plus important du corpus.
D'une part, les incidents ne remontent pas. Celui de J1 a été raconté à un collègue de même
niveau, sur le registre de la plaisanterie, et n'a jamais été porté à la connaissance de son
encadrante. Les notes consignent le motif~: la honte et la crainte d'un reproche
d'inattention. D'autre part, le contrôle s'adapte, mais individuellement et après coup. E1
contrôle désormais le processus et non plus le seul contenu final, demandant ce qui relève du
collaborateur, ce qui relève de l'outil et si les sources ont été vérifiées. E3 a formalisé
pour lui-même un contrôle en trois temps~: cohérence interne, contrôle croisé des sources,
test de défendabilité devant le client. E2, qui n'a pas vécu d'incident marquant, indique ne
pas contrôler différemment un contenu assisté. Aucune de ces adaptations n'est écrite,
partagée ni opposable.

\subsection{Les deux clivages}

Le premier a été rencontré axe par axe~: les juniors et les encadrants observent le même
phénomène et ne se le disent pas. J1 s'inquiète de son propre apprentissage, E1 de la
profondeur des livrables qu'il reçoit, et le corpus ne comporte aucune trace d'une
conversation entre ces deux positions. Le silence est lui-même documenté, dans des termes
presque identiques et sans concertation~: J1 note que tout le monde utilise l'outil et que
personne ne le dit, et rattache ce silence à la crainte du reproche de ne plus réfléchir. J2
emploie l'expression de «~secret de polichinelle~». Ce non-dit interne, codé en T3.1, n'était
pas anticipé. La grille prévoyait l'opacité vis-à-vis du client, non la dissimulation
entre collègues. Le résultat n'est pas un désaccord entre positions, c'est l'absence de
l'échange qui permettrait d'en avoir un.

Le second est plus inattendu~: les encadrants ne forment pas un bloc. E1, manager de dix ans,
n'utilise presque pas l'outil, se déclare plus rapide seul et emploie le terme de «~béquille~»
à propos des juniors. E2, associé de vingt-cinq ans et sponsor de la practice
concernée, l'utilise quotidiennement et le décrit comme un «~accélérateur de compétence, pas
substitut~». La posture ne suit donc ni l'ancienneté ni le grade, et l'écart se lit dans la
pratique de contrôle décrite ci-dessus, ce qu'enregistre le code T3.2.

Un dernier élément doit être signalé sans être tranché. Le chapitre~1 décrit un cadre
reposant sur un seul outil homologué. Or E2 précise ne jamais verser de données
confidentielles dans des outils non officiels, formulation qui suppose un recours à ces outils
pour les contenus qui ne sont pas confidentiels. J1, lui, déclare un usage officieux dont il doute que
la prudence qui l'accompagne soit générale. La règle n'est donc pas comprise de la même façon
selon la position occupée. Statuer sur la conformité de ces pratiques dépasse ce que le
dispositif permet d'établir. Le constat est autre~: un cadre diversement interprété
n'encadre pas.

\section{Confrontation aux cadres théoriques}

\subsection{Le seuil de spécificité, ou la frontière dentelée décrite de l'intérieur}

Dell'Acqua et ses coauteurs \cite{dellacqua2026} établissent, sur une population de
consultants, l'existence d'une frontière irrégulière séparant les tâches sur lesquelles
l'outil produit un gain net de celles sur lesquelles il dégrade la performance, et montrent
qu'elle est difficile à percevoir de l'intérieur. Le corpus décrit exactement cette
frontière, et la décrit cinq fois, par des personnes qui n'ont pas lu l'article.

Deux nuances font l'intérêt de la confrontation. Les répondants ont localisé la frontière, là
où l'expérimentation la donne pour peu visible. Mais la méthode est unique et coûteuse~:
l'échec chiffré en minutes perdues puis mémorisé en règle personnelle. La localisation est
acquise~; elle a été payée cinq fois, séparément, par des personnes qui ignoraient que leurs
collègues avaient payé le même prix. Ensuite, cette frontière ne sépare pas les tâches faciles
des tâches difficiles, mais celles dont le résultat se vérifie par recoupement de celles qui
exigent un contexte non observé. C'est la frontière de la sanctuarisation décrite en 4.1~:
les répondants n'ont donc pas construit deux règles mais une seule.

Un désalignement mérite enfin d'être noté. Là où Noy et Zhang \cite{noy2023} puis
\mbox{Brynjolfsson}, Li et Raymond \cite{brynjolfsson2025} montrent des gains concentrés sur les
moins expérimentés, le gain le plus élevé est ici déclaré par le répondant le plus ancien. La
variable discriminante est le portefeuille de tâches plutôt que le grade.

\subsection{Un contrôle qui s'adapte par l'incident et non par la règle}

Kellogg, Valentine et Christin \cite{kellogg2020} analysent le contrôle algorithmique comme un
terrain de lutte, en recensant les leviers par lesquels une technologie oriente et discipline
le travail et les résistances qu'ils suscitent. Le terrain observé se laisse nommer par défaut
dans cette grille~: il n'y a pas ici de contrôle algorithmique du travail, mais un travail
assisté sur lequel aucun dispositif de contrôle n'a été construit. Le cadre existant est un
cadre d'autorisation (tel outil est permis, tel autre non) et non un cadre de
vérification. Le contournement se produit donc en amont de toute règle contraignante~: le
non-dit codé en T3.1 est social et non procédural, puisqu'aucune procédure ne prescrit de
déclarer quoi que ce soit.

Lebovitz, Lifshitz-Assaf et Levina \cite{lebovitz2022} éclairent le versant individuel, en
montrant comment des professionnels arbitrent face à l'opacité d'un outil au moment de porter
un jugement qui les engage. Le protocole en trois temps d'E3 est exactement un dispositif de
ce type. La règle qu'il transmet aux juniors en est la maxime~:
«~si tu peux pas justifier, mets pas~». Ce dispositif existe donc, il est fin et efficace~;
mais il est privé,
construit par un individu à ses frais, et rien dans l'organisation ne le rend disponible aux
autres. La difficulté que soulèvent Lebovitz, Levina et Lifshitz-Assaf \cite{lebovitz2021} à
propos de l'étalon de vérification se retrouve sous forme pratique~: ce qui a sauvé J1 lors de
son quasi-incident est d'avoir conservé une source de vérité à côté de lui, circonstance
heureuse et non dispositif.

Le mécanisme d'ensemble peut alors être énoncé. Le circuit de revue du cabinet est hérité de
la standardisation des qualifications propre à la bureaucratie professionnelle
\cite{mintzberg1982}, où le contrôle porte sur la personne formée plutôt que sur l'acte~; il
tient tant que le relecteur sait quoi chercher. Or la nature des erreurs a changé (E1
relève que celles de l'outil sont plus subtiles que les erreurs humaines) sans que le
contenu de la qualification ait suivi. La compétence de vérification s'acquiert donc par
l'expérience, et cette expérience est celle de l'incident.

\subsection{La relation junior-encadrant ne s'appauvrit pas~: elle change d'objet}

Le résultat le plus solide du corpus est celui qui contredit l'hypothèse formulée en amont. La
grille avait inscrit en B3.2 l'appauvrissement de l'interaction entre junior et encadrant,
cohérent avec ce qu'établit Anthony \cite{anthony2021} sur des analystes juniors qui, faute
d'interroger le fonctionnement d'une technologie opaque, développent une expertise plus
superficielle. Le mécanisme attendu était direct~: si la machine produit le premier jet, la
boucle de correction par le senior perd son objet et le compagnonnage s'étiole.

Ce n'est pas ce que le terrain donne à voir, et les deux positions le contredisent
indépendamment. J1 a pourtant toutes les raisons de noircir le tableau, puisqu'il est
celui qui parle d'«~illusion de compétence~». Il décrit sa première proposition commerciale
comme un moment de formation fait de commentaires de titres et de structure, et observe que
ces commentaires portent aujourd'hui sur la solidité du raisonnement. E1 décrit
symétriquement ce que transmet une relecture (le jugement, la raison pour laquelle un
chiffre est suspect) et l'estime plus importante encore avec l'outil.

L'interprétation la plus économe passe par le modèle de conversion des connaissances
\cite{nonaka1997}. L'hypothèse d'appauvrissement supposait que la socialisation, mode de
conversion du tacite vers le tacite, avait pour support le premier jet du junior~: la machine
produisant ce jet, le support disparaissait. Le corpus indique qu'il n'a pas disparu mais
changé de niveau. La forme étant déléguée, la revue se reporte sur ce qu'elle n'atteignait
auparavant qu'après plusieurs itérations~: la justification, la pertinence, la source. La
socialisation ne se réduit pas~; elle porte sur un objet de rang supérieur. Et ce qu'E1 décrit
(énoncer ce qui rend un chiffre suspect, ce qui doit être vérifié et pourquoi) relève
moins de la socialisation que de l'externalisation, conversion du tacite vers l'explicite~:
l'outil oblige l'encadrant à mettre en mots un jugement jusque-là incorporé dans sa pratique.
Le gain de transmissibilité est réel, et il est involontaire.

La portée du résultat doit être bornée. La revue doit avoir lieu et porter sur le processus,
ce qu'E1 fait depuis son incident et qu'E2 déclare ne pas faire~: le bénéfice tient à la
pratique de l'encadrant plutôt qu'à la technologie, et il n'est pas généralisé dans le panel.
Il coexiste ensuite avec une perte en amont que le même E1 signale, le chemin de réflexion
initial n'étant plus parcouru. Cette perte affecte l'internalisation, conversion de
l'explicite vers le tacite qui suppose l'exécution. La relation monte d'un cran~; ce sur quoi elle
s'exerce a été appauvri en amont. Le corpus documente enfin des perceptions et non des
acquisitions.

C'est ici que la règle des trois encadrants prend son sens théorique. Exécuter manuellement
une fois avant d'accélérer est la condition d'internalisation. Sans exécution, le savoir
explicite déposé dans les méthodologies ne se convertit pas en savoir tacite. Le
professionnel se retrouve alors dans la situation, nommée en B4.2, de valider un contenu
qu'il n'aurait pas su produire. Or la standardisation des qualifications, mécanisme de
coordination principal d'une bureaucratie professionnelle \cite{mintzberg1982}, est dans le
conseil produite par le cabinet lui-même faute d'institution externe qui la délivre
\cite{nordenflycht2010}. Ce que la section 2.1.2 établissait sur le plan structurel se
vérifie donc ici~: la perturbation du parcours d'apprentissage n'atteint pas une fonction
support, elle atteint le mécanisme de coordination lui-même.

\subsection{Un point plus fragile~: la légitimité}

La confrontation à Abbott \cite{abbott1988} et à Suchman \cite{suchman1995} appelle plus de
réserve, faute de matériau. Les cinq répondants situent spontanément la valeur future du
cabinet là où Abbott place le noyau défendable d'une juridiction, dans l'inférence plutôt que
dans la collecte ou la mise en forme. Ceux qui refusent de mentionner l'outil au client
protègent un signal plutôt qu'une information, ce qui rejoint la lecture d'Alvesson
\cite{alvesson2001}. Mais cette convergence a été obtenue sur une question projective à dix
ans, dont la forme appelle une réponse valorisante. Les registres de Suchman permettent de
classer ce qui a été dit. Le registre cognitif est documenté par les cinq, chacun rapportant
la confiance du client à la fiabilité et à la maîtrise du domaine. Le moral l'est par trois,
par la responsabilité assumée devant le client, à quoi s'ajoute la non-déclaration de
l'usage, qui porte sur ce qui a été réellement fait et par qui. Le pragmatique n'est
documenté qu'une fois.
Ils ne permettent pas, en revanche, d'établir comment la légitimité se reconstruit devant
un client.

\section{Recommandations managériales}

Les recommandations qui suivent ne procèdent pas d'une opinion sur l'opportunité de ces
outils~: elles reprennent ce que le terrain a formulé, en le rendant opérationnel. Interrogés
séparément sur la décision qu'ils prendraient à la place de la direction, les cinq répondants
ont convergé sans s'être parlé et sur le même point d'équilibre. Aucun ne demande
d'interdiction, tous demandent que la vérification soit organisée, et E1 et E3 y ajoutent la
traçabilité des contenus produits avec assistance. Sur un
sujet où le clivage entre ceux qui produisent et ceux qui contrôlent était attendu, cette
convergence est elle-même un résultat.

\subsection{Séquencer l'autorisation d'usage dans le parcours du junior}

Chaque type de mission identifie une liste courte de tâches fondatrices (plan de livrable,
synthèse d'un corpus documentaire, chiffrage, note de cadrage) que le collaborateur produit
une fois sans assistance, la production étant relue et validée comme telle. Le passage est
acté~; l'usage de l'outil est ensuite non seulement autorisé mais encouragé sur ces mêmes
tâches, conformément à la séquence décrite par E2. La mesure doit se présenter comme une
séquence et non comme une restriction, faute de quoi elle produira les contournements
qu'anticipe la littérature sur les interventions organisationnelles \cite{venkatesh2008}.
Une règle vécue comme bridante sur une population déjà équipée est contournée en silence, et
le cabinet perd à la fois la formation et l'information.

\subsection{Rendre la vérification déclarative et graduer la relecture}

La mesure proposée par E1 est directement transposable~: le livrable produit avec assistance
porte une mention déclenchant une relecture renforcée. Elle répond à un constat saturé,
l'absence de toute étape de contrôle propre aux contenus générés. Trois conditions en
déterminent le succès. Le dispositif doit être déclaratif et non détectif~: il sert à
permettre le signalement de l'usage, non à le rechercher. Il doit être non punitif et énoncé
comme tel, faute de quoi il renforcera le non-dit interne au lieu de le lever. Le silence
actuel est motivé par la crainte du reproche, et un marquage perçu comme un aveu produirait
l'effet inverse. Il doit enfin être associé à la question qu'E1 pose déjà en revue, qui
distingue ce qui relève du collaborateur de ce qui relève de l'outil et s'assortit de la
vérification des sources. La finalité, telle qu'E3 la formule, est de responsabiliser et non
de contrôler.

\subsection{Constituer la vérification en compétence formée et outillée}

Quatre des cinq demandent une formation, et tous précisent qu'elle ne doit pas porter sur la
manipulation de l'outil. J1 la veut sur la détection des erreurs, le risque n'étant pas que
l'outil se trompe mais que le consultant ne le voie pas. E1 veut que les encadrants soient
formés à une relecture spécifique, les erreurs de l'outil étant plus subtiles. E2 la range du
côté de la compétence et non de la défiance~; E3 la veut obligatoire. Le cabinet n'a pas à en
concevoir le contenu, lequel existe déjà sous forme privée dans le corpus. Le protocole en
trois temps d'E3 est un standard éprouvé et transmissible en une page, que complète la
pratique de J1 consistant à conserver une source de vérité indépendante à côté du document en
production. La mesure consiste à expliciter ce savoir tacite~: rassembler les protocoles
individuels, les formaliser, les enseigner. La charte d'une page demandée par J2, qui
distingue ce qui est possible, ce qui ne l'est pas et ce qui doit être vérifié, en constitue
le format~; le canal d'échange que J2 et E2 réclament l'un et l'autre en constitue le
support d'entretien.

\subsection{Faire évoluer le dispositif autrement que par l'accident}

Le principe directeur qui commande les mesures précédentes s'énonce simplement~: un dispositif
de contrôle qui n'évolue que par accident est un dispositif qui attend l'accident. Le corpus
l'établit deux fois~: les encadrants qui ont modifié leur revue l'ont fait après un incident
vécu, et le seul quasi-incident rapporté par un junior n'est jamais remonté. Deux mesures en
découlent. La première est un traitement organisé des quasi-incidents, dissocié de toute
conséquence individuelle. Un incident détecté avant le client est une information de valeur,
et sa remontée doit être plus avantageuse que son silence. Les cas recueillis alimentent la
révision périodique du protocole de vérification, de sorte que l'apprentissage cesse de
dépendre de la répartition aléatoire des incidents. La seconde est la clarification du
périmètre normatif, dont le flou est signalé par les deux positions et produit un usage dont
l'ampleur échappe au cabinet. Cette évolution est enfin cohérente avec le repositionnement que
les cinq décrivent déjà (la valeur vendue se déplace vers le jugement et la responsabilité
assumée \cite{suchman1995}), car un dispositif de vérification explicite est ce qui rend
cette promesse vérifiable.

\section{Limites des résultats}

Les limites du dispositif ont été énoncées au chapitre précédent~; celles qui suivent portent
sur les résultats eux-mêmes.

La première tient à l'assise empirique, décrite en 3.4 et non reprise ici. Sa conséquence
sur les résultats est en revanche précise~: les convergences relevées ne valent pas comme
fréquences. Un constat dit saturé signifie que les cinq récits le documentent, non qu'il
est répandu dans la practice. Un constat dit clivant oppose des répondants nommément
identifiés, non deux groupes constitués, et le clivage sur l'adaptation de la revue passe à
l'intérieur du groupe des encadrants.

La deuxième tient à l'inégale solidité des axes. L'axe C est le plus mince et porte pourtant
la troisième question directrice. Un seul répondant rapporte une interpellation réelle par un
client, un seul une conséquence commerciale documentée. Le reste est du déclaratif projectif,
recueilli sur des questions dont la forme appelle une réponse valorisante. La convergence
obtenue sur cet axe est donc d'une autre nature que celle obtenue sur les axes A et D, où elle
repose sur des récits d'incidents datés et chiffrés. Elle dit ce que le métier devrait être,
non comment la légitimité se reconstruit devant un client.

La troisième tient au matériau. Les extraits reproduits ici sont des fragments de notes prises
en séance, donnés tels qu'ils ont été consignés ou rapportés au style indirect~: aucun n'a
valeur de citation littérale, et l'affectation d'un code dépend d'une formulation en partie
celle de l'enquêteur.

La quatrième tient à la position de l'enquêteur, décrite en 3.4, et affecte spécifiquement les deux clivages.
Cette position a probablement rendu possible le recueil
du non-dit codé en T3.1~: un enquêteur extérieur n'aurait guère obtenu de deux juniors
l'aveu d'un usage dissimulé à leur hiérarchie. Mais elle expose au risque symétrique de
reconnaître d'autant plus aisément un phénomène qu'on s'attendait à trouver. Le silence entre
positions est établi par ce que les répondants en disent, non par l'observation de leurs
échanges.

La cinquième porte sur le cadre normatif~: le dispositif ne permet ni d'établir l'ampleur des
usages relevés en 4.1, ni de statuer sur leur conformité. Ce point est versé au dossier comme
une observation de gouvernance appelant vérification, non comme un constat de manquement.

% ==============================================================================
% Fin du chapitre 4
% ==============================================================================

\cleardoublepage{}

% ==============================================================================
% PARTIE 3 : ANALYSE ET RECOMMANDATIONS
% ==============================================================================
\part{Analyse des Pratiques et Recommandations Managériales}

\chapter{Analyse des résultats : L'IA au quotidien, entre opportunisme et vigilance}
%\input{chapters/chapitre5.tex}

\chapter{Recommandations stratégiques et RH pour le cabinet}
%\input{chapters/chapitre6.tex}

\cleardoublepage{}

% ==============================================================================
% CONCLUSION GENERALE
% ==============================================================================
\chapter{Conclusion Générale}
% À rédiger en fin de parcours

\cleardoublepage{}

% ==============================================================================
% BIBLIOGRAPHIE (Gérée proprement via un fichier .bib séparé)
% ==============================================================================
\renewcommand{\tocbibname}{Bibliographie / Webographie}
\bibliography{references} % Crée un fichier references.bib dans ton projet
\bibliographystyle{plain}

\cleardoublepage{}

% Index des figures et tableaux (Optionnel mais requis si tu en ajoutes)
\listoffigures
\cleardoublepage{}

\listoftables
\cleardoublepage{}

% ==============================================================================
% ANNEXES & RÉSUMÉS DE FIN
% ==============================================================================
\appendix
\part*{Annexes} 
\addcontentsline{toc}{part}{Annexes} 
\cleardoublepage{}

\chapter{Guide d'entretien semi-directif}
%% ==============================================================================
% ANNEXE 1 -- GUIDE D'ENTRETIEN SEMI-DIRECTIF
% Mémoire MAE -- Arthur SAUNIER -- Wavestone Luxembourg
% Le \chapter{} est déclaré dans le fichier maître ; ce fichier démarre en \section.
% ==============================================================================

\section{Objectifs et architecture du guide}

Le présent guide constitue l'instrument de collecte de la phase empirique. Il a été
construit de manière déductive à partir des quatre questions directrices dérivées de
la question centrale de recherche, chaque thème d'entretien constituant la
traduction conversationnelle d'une de ces questions. Ce chaînage explicite entre
questionnement théorique et questionnement de terrain garantit la validité interne
du dispositif~: aucune question posée aux répondants n'est orpheline d'une
interrogation de recherche, et aucune question de recherche n'est laissée sans
matériau empirique.

L'entretien est calibré pour une durée de trente-cinq à quarante-cinq minutes,
compatible avec les contraintes de disponibilité de consultants en mission. Cette
contrainte impose une discipline de conduite~: les thèmes~3 et~4 étant les plus
exposés au risque d'être sacrifiés en fin d'entretien, l'enquêteur veille à borner
strictement les deux premiers.

\begin{table}[h]
\centering
\small
\begin{tabular}{|p{0.08\linewidth}|p{0.40\linewidth}|p{0.23\linewidth}|p{0.09\linewidth}|}
\hline
\textbf{Thème} & \textbf{Question directrice adossée} & \textbf{Cadre théorique mobilisé} & \textbf{Durée} \\
\hline
T0 & Cadrage du répondant (hors questionnement) & --- & 5 min \\
\hline
T1 & QD1 -- Appropriation réelle des outils d'IA générative dans les tâches d'analyse & Sociologie des usages, TAM (Davis, 1989) & 10 min \\
\hline
T2 & QD2 -- Impact sur l'apprentissage par la pratique et la montée en compétences & Nonaka \& Takeuchi (1995), Polanyi (1966) & 10 min \\
\hline
T3 & QD3 -- Reconstruction de la légitimité de l'expert face au client & Abbott (1988), Suchman (1995) & 8 min \\
\hline
T4 & QD4 -- Encadrement managérial des risques sans bridage de l'autonomie & Mintzberg (1982), gouvernance des risques & 8 min \\
\hline
T5 & Clôture projective et ouverture & --- & 4 min \\
\hline
\end{tabular}
\caption{Architecture du guide d'entretien et adossement théorique}
\label{tab:architecture-guide}
\end{table}

\section{Consignes de conduite adressées à l'enquêteur}

La qualité du matériau recueilli dépend moins de la formulation des questions que de
la discipline de l'enquêteur. Les principes suivants encadrent la conduite de chaque
entretien.

\begin{itemize}
    \item \textbf{Privilégier le récit à l'opinion.} Chaque thème s'ouvre par une
    question de type narratif ou par un incident critique (« racontez-moi la
    dernière fois que… ») plutôt que par une question d'opinion. Le récit d'une
    situation vécue produit un matériau vérifiable et échappe aux réponses de
    façade que suscite un sujet socialement connoté.

    \item \textbf{Ne jamais introduire le vocabulaire de la recherche.} Les termes
    « légitimité », « montée en compétences », « hallucination », « gouvernance »
    ne doivent pas être prononcés par l'enquêteur avant que le répondant ne les ait
    lui-même mobilisés. Leur emploi prématuré induit la réponse et détruit la valeur
    probante du verbatim.

    \item \textbf{Tenir une position de neutralité axiologique.} L'enquêteur, lui-même
    consultant du cabinet, s'abstient de défendre comme de critiquer l'usage de ces
    outils. Toute marque d'approbation ou de réprobation oriente durablement la suite
    de l'entretien.

    \item \textbf{Exploiter le silence.} Une pause de trois à cinq secondes après une
    réponse jugée terminée produit fréquemment la relance la plus riche de l'entretien.

    \item \textbf{Relancer par la reformulation.} Les relances privilégiées sont
    l'écho (reprise du dernier segment de phrase sur un ton interrogatif), la demande
    de précision factuelle (« sur quelle mission~? ») et la demande d'exemple
    contraire (« et une fois où cela ne s'est pas passé ainsi~? »).

    \item \textbf{Assumer la souplesse de l'ordre.} L'ordre des thèmes est indicatif.
    Si le répondant aborde spontanément un thème ultérieur, l'enquêteur le suit et
    coche mentalement le thème traité plutôt que de le ramener au fil prévu.

    \item \textbf{Ne pas combler les contradictions.} Les incohérences entre le
    discours déclaré et les pratiques décrites constituent le matériau analytique le
    plus précieux~; elles sont notées, non résolues en séance.
\end{itemize}

\section{Préambule, anonymisation et recueil du consentement}

Le texte suivant est lu au répondant avant tout enregistrement.

\begin{quote}
\textit{Merci d'avoir accepté cet entretien. Il s'inscrit dans le cadre d'un mémoire
de Master~2 en Management et Administration des Entreprises, réalisé à l'IAE Nancy
parallèlement au stage effectué dans le cabinet. Ce travail porte sur la manière dont
les consultants intègrent les outils d'intelligence artificielle générative dans leur
travail d'analyse, et sur ce que cela change dans l'exercice du métier.}

\textit{Trois précisions avant de commencer. Premièrement, il ne s'agit ni d'un audit,
ni d'une évaluation~: il n'existe pas de bonne ou de mauvaise réponse, et la
description de pratiques informelles est aussi utile que celle des usages officiels.
Deuxièmement, l'entretien est intégralement anonymisé~: aucun nom de personne, de
client ou de mission ne figurera dans le mémoire, et les répondants y sont désignés
par un code de type « consultant junior~1 ». Troisièmement, l'entretien dure environ
quarante minutes et l'enregistrement audio, s'il est autorisé, sert uniquement à la
retranscription et est détruit après l'analyse.}

\textit{Acceptez-vous que l'entretien soit enregistré~?}
\end{quote}

En cas de refus d'enregistrement, l'entretien se poursuit sous prise de notes
manuscrite, la restitution des verbatims étant alors signalée comme reconstituée.

\section{Thème 0 --- Trajectoire et cadre d'exercice}

Ce thème ne relève d'aucune question directrice. Il vise à installer la parole, à
recueillir les variables de contrôle nécessaires à la lecture croisée des profils et à
identifier le vocabulaire propre au répondant, que l'enquêteur réemploiera ensuite.

\begin{itemize}
    \item Pour commencer, pouvez-vous décrire votre parcours et la manière dont vous
    êtes arrivé dans la practice~?
    \item Quel type de missions occupe l'essentiel de votre temps en ce moment~?
    \item Sur une mission typique, en quoi consiste concrètement votre journée~?
    \item Quelle part de votre travail relève de la production de livrables écrits~?
\end{itemize}

\section{Thème 1 --- L'usage réel des outils dans le travail d'analyse}

\textit{Question directrice~1~: de quelle manière les consultants s'approprient-ils
les outils d'intelligence artificielle générative dans la réalisation de leurs tâches
d'analyse quotidiennes~?}

\subsection*{Question d'ouverture}

Pouvez-vous me raconter la dernière fois que vous avez utilisé un de ces outils dans
votre travail~? Que cherchiez-vous à faire, et comment cela s'est-il passé~?

\subsection*{Relances de premier niveau}

\begin{itemize}
    \item Sur quelles tâches précisément ces outils interviennent-ils dans votre
    quotidien~?
    \item Comment avez-vous appris à vous en servir~? Par une formation, par un
    collègue, seul~?
    \item À quel moment de la production d'un livrable les mobilisez-vous~: au
    démarrage, en cours de rédaction, à la relecture~?
    \item Y a-t-il des tâches pour lesquelles vous ne les utilisez jamais~? Lesquelles,
    et pourquoi~?
    \item Utilisez-vous les outils mis à disposition par le cabinet, ou d'autres~?
    Qu'est-ce qui fait la différence pour vous~?
\end{itemize}

\subsection*{Relances de second niveau}

À n'utiliser que si le répondant demeure dans la généralité.

\begin{itemize}
    \item Vous parlez de gain de temps~: sur la dernière mission, à quoi avez-vous
    consacré le temps ainsi libéré~?
    \item Y a-t-il des situations où cela vous a demandé plus d'effort que de le faire
    vous-même~?
    \item Vous êtes-vous déjà arrêté en cours d'utilisation en vous disant que ce
    n'était pas la bonne méthode~? Qu'est-ce qui a déclenché ce moment~?
    \item En parlez-vous avec vos collègues~? Est-ce un sujet dont on discute
    ouvertement dans l'équipe~?
\end{itemize}

\section{Thème 2 --- Apprentissage du métier et construction de l'expertise}

\textit{Question directrice~2~: quel est l'impact de cette automatisation sur
l'apprentissage par la pratique et la trajectoire de montée en compétences des
consultants juniors~?}

\subsection*{Question d'ouverture}

Si vous repensez à la manière dont vous avez appris ce métier --- ou dont vous voyez
les nouveaux arrivants l'apprendre aujourd'hui ---, qu'est-ce qui a changé~?

\subsection*{Relances de premier niveau}

\begin{itemize}
    \item Quelles sont, selon vous, les tâches par lesquelles on apprend réellement le
    métier de consultant~?
    \item Ces tâches sont-elles encore réalisées de la même manière aujourd'hui~?
    \item Comment vous y prenez-vous pour vérifier ce que produit l'outil~? Que faut-il
    savoir pour être capable de le vérifier~?
    \item Un consultant qui débuterait aujourd'hui avec ces outils apprendrait-il les
    mêmes choses que vous~?
    \item Comment se passe la relecture d'un livrable par un senior~: qu'est-ce qui se
    transmet à ce moment-là~?
\end{itemize}

\subsection*{Relances de second niveau}

\begin{itemize}
    \item Vous êtes-vous déjà retrouvé à valider un contenu que vous n'auriez pas su
    produire vous-même~? Comment avez-vous procédé~?
    \item Y a-t-il des compétences que vous sentez que vous exercez moins qu'avant~?
    \item Si l'accès à ces outils était coupé demain sur une mission, qu'est-ce qui
    deviendrait difficile pour l'équipe~?
\end{itemize}

\section{Thème 3 --- La relation au client et la valeur de l'expertise}

\textit{Question directrice~3~: comment l'autorité et la légitimité de l'expert humain
se reconstruisent-elles vis-à-vis du client lorsque la machine produit une partie
substantielle de la valeur intellectuelle~?}

\subsection*{Question d'ouverture}

Quand vous présentez un livrable à un client, sur quoi repose selon vous sa confiance
dans ce que vous lui remettez~?

\subsection*{Relances de premier niveau}

\begin{itemize}
    \item Le sujet de l'usage de ces outils a-t-il déjà été abordé avec un client~?
    Dans quelles circonstances, et comment cela s'est-il passé~?
    \item Diriez-vous à un client qu'une partie de l'analyse a été produite avec ces
    outils~? Qu'est-ce qui guiderait cette décision~?
    \item Est-ce que les clients eux-mêmes utilisent ces outils~? Cela change-t-il
    quelque chose à ce qu'ils attendent du cabinet~?
    \item Qu'est-ce que le cabinet vend aujourd'hui que le client ne pourrait pas
    obtenir autrement~?
\end{itemize}

\subsection*{Relances de second niveau}

\begin{itemize}
    \item Avez-vous déjà eu le sentiment qu'un livrable était perçu comme moins
    crédible~? Qu'est-ce qui l'a provoqué~?
    \item Dans dix ans, sur quoi un client acceptera-t-il encore de payer un cabinet
    de conseil~?
\end{itemize}

\section{Thème 4 --- Encadrement, contrôle qualité et gestion des risques}

\textit{Question directrice~4~: comment le management encadre-t-il l'usage de l'IA
pour pallier les risques opérationnels sans brider l'autonomie des collaborateurs~?}

\subsection*{Question d'ouverture}

Existe-t-il des règles, écrites ou non, sur l'usage de ces outils dans le cabinet~?
Comment en avez-vous eu connaissance~?

\subsection*{Relances de premier niveau}

\begin{itemize}
    \item Que se passe-t-il concrètement avant qu'un livrable ne parte chez le client~?
    Qui contrôle quoi~?
    \item Avez-vous déjà rencontré un contenu produit par un outil qui s'est révélé
    faux~? Que s'est-il passé ensuite~?
    \item Comment traitez-vous la question des données de mission lorsque vous
    utilisez ces outils~?
    \item Ces règles vous paraissent-elles adaptées au travail réel~? Y a-t-il des
    situations où elles sont contournées~?
    \item Selon vous, qui est responsable si une erreur passe dans un livrable remis
    au client~?
\end{itemize}

\subsection*{Relances de second niveau}

\begin{itemize}
    \item Si vous aviez à écrire vous-même la règle, que diriez-vous~?
    \item Y a-t-il un équilibre à trouver entre encadrer et laisser faire~? Où le
    placeriez-vous~?
\end{itemize}

\section{Thème 5 --- Clôture}

\begin{itemize}
    \item Si vous étiez à la place de la direction du cabinet, quelle serait la
    première décision que vous prendriez sur ce sujet~?
    \item Y a-t-il un aspect important que nous n'avons pas abordé et que vous
    souhaiteriez ajouter~?
    \item Voyez-vous une personne du cabinet dont le point de vue serait
    particulièrement utile à recueillir~?
\end{itemize}

L'entretien se clôt par le rappel des modalités d'anonymisation et par la proposition
d'une transmission de la synthèse finale au répondant.

\section{Variantes de formulation selon le profil}

Le guide demeure identique dans sa structure quel que soit le répondant, condition de
la comparabilité des matériaux. Seules les formulations d'ancrage sont adaptées afin
de solliciter la position d'énonciation propre à chaque profil.

Le panel distingue deux positions, et cette distinction épouse celle du dispositif de
revue décrit au chapitre~1. Les \textbf{juniors} — analystes et consultants — rédigent les
livrables. Les \textbf{encadrants} — consultants seniors, manager et senior manager —
accompagnent la production puis la contrôlent avant l'envoi au client. Le contraste
recherché est celui de deux positions dans la chaîne de production, non celui de deux
échelons hiérarchiques.

\begin{table}[h]
\centering
\small
\begin{tabular}{|p{0.08\linewidth}|p{0.38\linewidth}|p{0.38\linewidth}|}
\hline
\textbf{Thème} & \textbf{Juniors (analystes, consultants)} & \textbf{Encadrants (consultants seniors, manager, senior manager)} \\
\hline
T1 & Usage vécu et découverte des outils & Usage propre \emph{et} usages observés dans l'équipe \\
\hline
T2 & Sentiment sur son propre apprentissage & Perception de l'évolution des juniors encadrés \\
\hline
T3 & Position en soutenance de livrable & Position en relation client directe \\
\hline
T4 & Connaissance et vécu des règles & Rôle effectif dans le contrôle qualité, arbitrage entre risque et productivité \\
\hline
\end{tabular}
\caption{Adaptation des angles d'interrogation selon la position du répondant}
\label{tab:variantes-profils}
\end{table}

\section{Fiche de traçabilité de l'entretien}

Une fiche est renseignée immédiatement après chaque entretien, avant toute
retranscription, afin de conserver les éléments contextuels non enregistrés.

\begin{itemize}
    \item Code du répondant, profil et ancienneté approximative dans le conseil.
    \item Date, durée effective, modalité (présentiel ou visioconférence).
    \item Autorisation d'enregistrement accordée ou refusée.
    \item Thèmes effectivement couverts et thèmes écourtés.
    \item Climat de l'entretien et éventuelles réticences perçues.
    \item Trois formulations marquantes notées à chaud.
    \item Hypothèses émergentes à tester lors des entretiens suivants.
\end{itemize}

\cleardoublepage{}
\thispagestyle{empty}

% Résumés réglementaires de quatrième de couverture (Consigne IAE)
\section*{Résumé}
\addcontentsline{toc}{chapter}{Résumé}
Ce mémoire professionnel s'intéresse à l'impact de l'intégration des outils d'intelligence artificielle générative au sein des cabinets de conseil... [À compléter en fin de rédaction]

\vspace{1cm}
{\bf Mots-clés~:} Intelligence Artificielle Générative, Sociétés de Conseil, Gestion des Compétences, Bureaucratie Professionnelle, Légitimité.

\section*{Abstract}
\addcontentsline{toc}{chapter}{Abstract}
This professional thesis explores the impact of integrating generative artificial intelligence tools within management and IT consulting firms... [À compléter en fin de rédaction]

\vspace{1cm}
{\bf Keywords~:} Generative Artificial Intelligence, Management Consulting, Skills Management, Professional Bureaucracy, Legitimacy.

\printSignatures

\end{document}